\documentclass[../main.tex]{subfiles}
\begin{document}

\chapter{Memory Management}
\lessoninfo{
  video={P3L2},
  textbook={OSTEP \cite{ostep} ch.~13, 15--23; Silberschatz et al.\ \cite{silberschatz2018} ch.~9--10; Bovet and Cesati \cite{bovet2005} ch.~2, 8, 17},
  keywords={virtual memory, paging, segmentation, page table entry, page fault, multi-level page table, TLB, large page, buddy allocator, slab allocator, demand paging, page replacement, copy-on-write, checkpointing},
}

Lesson~7 described how the operating system shares the CPU among processes.
This lesson turns to the second resource that every process needs, physical
memory.  The operating system gives every process its own virtual address
space (Lesson~2), keeps in physical memory only the parts that are in use,
and brings other parts in from secondary storage when they are needed.  We
examine the two ways of organizing this mapping, pages and segments; the
hardware and the data structures that make address translation fast and
compact; the allocators that manage physical memory; the policies that
decide which pages leave memory; and two services, copy-on-write and
checkpointing, that the memory management hardware makes possible.

\section{Goals of memory management}
\label{sec:l08-goals}

The operating system manages the physical memory (DRAM) of the machine on
behalf of the processes that run on it.\source{P3L2, lesson preview,
memory management: goals; OSTEP \cite{ostep} ch.~13, 15; Silberschatz et
al.\ \cite{silberschatz2018} ch.~9.}  Processes do not use physical
addresses: each has a virtual address space whose range can be far larger
than physical memory.  With 32-bit virtual addresses every process can
address \SI{4}{\gibi\byte}, and the address spaces of all processes
together easily exceed the memory that is installed.

\begin{definition}[Virtual memory]
  The abstraction with which the operating system decouples the virtual
  address spaces of processes from physical memory is called
  \term{virtual memory}[仮想メモリ].  The operating system maps the parts
  of every address space that are in use onto physical memory, keeps other
  parts on secondary storage, and translates and validates every address
  that a process uses.
\end{definition}

A process does not need all of its state at once; at any time it uses only
part of its address space.  The operating system can therefore keep only
that part in physical memory, as long as the state that the process needs
next can be made available quickly.  Two tasks follow.
\begin{description}
  \item[Allocation] The operating system decides which portions of which
    address spaces reside in physical memory and where, and it establishes
    the mapping from virtual to physical addresses.  Since physical memory
    may be smaller than the memory the processes use, some state must be
    moved to disk and brought back when it is needed; allocation therefore
    includes \emph{replacement}, the choice of the state that leaves
    memory to make room.
  \item[Arbitration] Every memory access must be translated from a virtual
    into a physical address, and validated: the address must belong to an
    allocated region of the address space, and the access must be
    permitted.  The conversion of a virtual address into the physical
    address to which it maps is called
    \term{address translation}[アドレス変換].  It happens on every memory
    access and must therefore be fast.
\end{description}

The operating system can divide the address space into units in two ways
(\cref{fig:l08-goals}).
\begin{description}
  \item[Page-based memory management] With \term{paging}[ページング] the
    virtual address space is divided into fixed-size blocks, each called a
    \term{page}[ページ], and physical memory into blocks of the same size,
    each called a \term{page frame}[ページフレーム].  Allocation maps pages
    to page frames; arbitration uses page tables.
  \item[Segment-based memory management] With
    \term{segmentation}[セグメンテーション] the address space is divided
    into blocks of flexible size, each called a
    \term{segment}[セグメント], which correspond to logical components of
    the process, such as its code, heap or stack.  Allocation maps segments
    to regions of physical memory of the same size; arbitration uses
    segment registers.
\end{description}
Paging is the dominant approach in current operating systems.  Segmentation
is supported on some architectures, notably x86, where it is combined with
paging (\cref{sec:l08-segmentation}).

\begin{figure}[htbp]
  \centering
  % Page-based versus segment-based memory management: fixed-size pages
% mapped to page frames (one page on disk, one unmapped) and variable-size
% segments mapped to regions of physical memory.
% Usage: % Page-based versus segment-based memory management: fixed-size pages
% mapped to page frames (one page on disk, one unmapped) and variable-size
% segments mapped to regions of physical memory.
% Usage: % Page-based versus segment-based memory management: fixed-size pages
% mapped to page frames (one page on disk, one unmapped) and variable-size
% segments mapped to regions of physical memory.
% Usage: \input{figures/l08-goals}   (width ~118 mm)
\begin{tikzpicture}[x=1mm, y=1mm, font=\scriptsize\sffamily,
  >={Stealth[length=1.4mm]},
  pg/.style={draw, fill=ln-accent!22},
  sg/.style={draw, fill=ln-source!16},
  oth/.style={draw, fill=ln-box},
  lab/.style={font=\scriptsize\sffamily},
  hd/.style={font=\scriptsize\sffamily\bfseries, anchor=base},
  capt/.style={font=\footnotesize\sffamily}]
  % ================= (a) paging
  \node[hd] at (9,38) {\iflangja{仮想}{virtual}};
  \node[hd] at (45,38) {\iflangja{物理}{physical}};
  % virtual pages 0..5 (height 6)
  \foreach \j/\s in {0/pg, 1/pg, 2/pg, 4/pg, 5/pg} {
    \draw[\s] (2,6*\j) rectangle ++(14,6);
    \node[lab] at (9,3+6*\j) {\iflangja{ページ}{page} \j};
  }
  \draw[fill=white] (2,18) rectangle ++(14,6);
  \node[lab, text=ln-muted] at (9,21) {\iflangja{未使用}{unused}};
  % page frames 0..7 (height 4.5)
  \foreach \i in {0,...,7} { \draw[oth] (38,4.5*\i) rectangle ++(14,4.5); }
  \foreach \i in {0,2,5,7} { \draw[pg] (38,4.5*\i) rectangle ++(14,4.5); }
  \foreach \i in {1,3,4,6} { \node[lab, text=ln-muted] at (45,2.25+4.5*\i) {\iflangja{他}{other}}; }
  % disk
  \draw[fill=ln-box, rounded corners=2pt] (38,-11) rectangle (52,-4)
    node[midway, lab] {\iflangja{ディスク}{disk}};
  % mappings: page -> frame
  \foreach \j/\i in {0/5, 1/2, 4/7, 5/0} {
    \draw[->] (16,3+6*\j) -- (38,2.25+4.5*\i);
  }
  \draw[->] (16,15) -- (38,-7.5);
  \node[capt] at (27,-15) {\iflangja{(a) ページング}{(a) paging}};
  % ================= (b) segmentation
  \node[hd] at (72,38) {\iflangja{仮想}{virtual}};
  \node[hd] at (108,38) {\iflangja{物理}{physical}};
  \draw[sg] (65,0)  rectangle (79,9)  node[midway, lab] {\iflangja{コード}{code}};
  \draw[sg] (65,9)  rectangle (79,14) node[midway, lab] {\iflangja{データ}{data}};
  \draw[sg] (65,14) rectangle (79,22) node[midway, lab] {\iflangja{ヒープ}{heap}};
  \draw[fill=white] (65,22) rectangle (79,30);
  \draw[sg] (65,30) rectangle (79,36) node[midway, lab] {\iflangja{スタック}{stack}};
  % physical memory
  \draw[oth] (101,-11) rectangle (115,36);
  \draw[sg] (101,-9) rectangle (115,-3) node[midway, lab] {\iflangja{スタック}{stack}};
  \draw[sg] (101,2)  rectangle (115,11) node[midway, lab] {\iflangja{コード}{code}};
  \draw[sg] (101,16) rectangle (115,24) node[midway, lab] {\iflangja{ヒープ}{heap}};
  \draw[sg] (101,28) rectangle (115,33) node[midway, lab] {\iflangja{データ}{data}};
  \draw[->] (79,4.5) -- (101,6.5);
  \draw[->] (79,11.5) -- (101,30.5);
  \draw[->] (79,18) -- (101,20);
  \draw[->] (79,33) -- (101,-6);
  \node[capt] at (90,-15) {\iflangja{(b) セグメンテーション}{(b) segmentation}};
\end{tikzpicture}
   (width ~118 mm)
\begin{tikzpicture}[x=1mm, y=1mm, font=\scriptsize\sffamily,
  >={Stealth[length=1.4mm]},
  pg/.style={draw, fill=ln-accent!22},
  sg/.style={draw, fill=ln-source!16},
  oth/.style={draw, fill=ln-box},
  lab/.style={font=\scriptsize\sffamily},
  hd/.style={font=\scriptsize\sffamily\bfseries, anchor=base},
  capt/.style={font=\footnotesize\sffamily}]
  % ================= (a) paging
  \node[hd] at (9,38) {\iflangja{仮想}{virtual}};
  \node[hd] at (45,38) {\iflangja{物理}{physical}};
  % virtual pages 0..5 (height 6)
  \foreach \j/\s in {0/pg, 1/pg, 2/pg, 4/pg, 5/pg} {
    \draw[\s] (2,6*\j) rectangle ++(14,6);
    \node[lab] at (9,3+6*\j) {\iflangja{ページ}{page} \j};
  }
  \draw[fill=white] (2,18) rectangle ++(14,6);
  \node[lab, text=ln-muted] at (9,21) {\iflangja{未使用}{unused}};
  % page frames 0..7 (height 4.5)
  \foreach \i in {0,...,7} { \draw[oth] (38,4.5*\i) rectangle ++(14,4.5); }
  \foreach \i in {0,2,5,7} { \draw[pg] (38,4.5*\i) rectangle ++(14,4.5); }
  \foreach \i in {1,3,4,6} { \node[lab, text=ln-muted] at (45,2.25+4.5*\i) {\iflangja{他}{other}}; }
  % disk
  \draw[fill=ln-box, rounded corners=2pt] (38,-11) rectangle (52,-4)
    node[midway, lab] {\iflangja{ディスク}{disk}};
  % mappings: page -> frame
  \foreach \j/\i in {0/5, 1/2, 4/7, 5/0} {
    \draw[->] (16,3+6*\j) -- (38,2.25+4.5*\i);
  }
  \draw[->] (16,15) -- (38,-7.5);
  \node[capt] at (27,-15) {\iflangja{(a) ページング}{(a) paging}};
  % ================= (b) segmentation
  \node[hd] at (72,38) {\iflangja{仮想}{virtual}};
  \node[hd] at (108,38) {\iflangja{物理}{physical}};
  \draw[sg] (65,0)  rectangle (79,9)  node[midway, lab] {\iflangja{コード}{code}};
  \draw[sg] (65,9)  rectangle (79,14) node[midway, lab] {\iflangja{データ}{data}};
  \draw[sg] (65,14) rectangle (79,22) node[midway, lab] {\iflangja{ヒープ}{heap}};
  \draw[fill=white] (65,22) rectangle (79,30);
  \draw[sg] (65,30) rectangle (79,36) node[midway, lab] {\iflangja{スタック}{stack}};
  % physical memory
  \draw[oth] (101,-11) rectangle (115,36);
  \draw[sg] (101,-9) rectangle (115,-3) node[midway, lab] {\iflangja{スタック}{stack}};
  \draw[sg] (101,2)  rectangle (115,11) node[midway, lab] {\iflangja{コード}{code}};
  \draw[sg] (101,16) rectangle (115,24) node[midway, lab] {\iflangja{ヒープ}{heap}};
  \draw[sg] (101,28) rectangle (115,33) node[midway, lab] {\iflangja{データ}{data}};
  \draw[->] (79,4.5) -- (101,6.5);
  \draw[->] (79,11.5) -- (101,30.5);
  \draw[->] (79,18) -- (101,20);
  \draw[->] (79,33) -- (101,-6);
  \node[capt] at (90,-15) {\iflangja{(b) セグメンテーション}{(b) segmentation}};
\end{tikzpicture}
   (width ~118 mm)
\begin{tikzpicture}[x=1mm, y=1mm, font=\scriptsize\sffamily,
  >={Stealth[length=1.4mm]},
  pg/.style={draw, fill=ln-accent!22},
  sg/.style={draw, fill=ln-source!16},
  oth/.style={draw, fill=ln-box},
  lab/.style={font=\scriptsize\sffamily},
  hd/.style={font=\scriptsize\sffamily\bfseries, anchor=base},
  capt/.style={font=\footnotesize\sffamily}]
  % ================= (a) paging
  \node[hd] at (9,38) {\iflangja{仮想}{virtual}};
  \node[hd] at (45,38) {\iflangja{物理}{physical}};
  % virtual pages 0..5 (height 6)
  \foreach \j/\s in {0/pg, 1/pg, 2/pg, 4/pg, 5/pg} {
    \draw[\s] (2,6*\j) rectangle ++(14,6);
    \node[lab] at (9,3+6*\j) {\iflangja{ページ}{page} \j};
  }
  \draw[fill=white] (2,18) rectangle ++(14,6);
  \node[lab, text=ln-muted] at (9,21) {\iflangja{未使用}{unused}};
  % page frames 0..7 (height 4.5)
  \foreach \i in {0,...,7} { \draw[oth] (38,4.5*\i) rectangle ++(14,4.5); }
  \foreach \i in {0,2,5,7} { \draw[pg] (38,4.5*\i) rectangle ++(14,4.5); }
  \foreach \i in {1,3,4,6} { \node[lab, text=ln-muted] at (45,2.25+4.5*\i) {\iflangja{他}{other}}; }
  % disk
  \draw[fill=ln-box, rounded corners=2pt] (38,-11) rectangle (52,-4)
    node[midway, lab] {\iflangja{ディスク}{disk}};
  % mappings: page -> frame
  \foreach \j/\i in {0/5, 1/2, 4/7, 5/0} {
    \draw[->] (16,3+6*\j) -- (38,2.25+4.5*\i);
  }
  \draw[->] (16,15) -- (38,-7.5);
  \node[capt] at (27,-15) {\iflangja{(a) ページング}{(a) paging}};
  % ================= (b) segmentation
  \node[hd] at (72,38) {\iflangja{仮想}{virtual}};
  \node[hd] at (108,38) {\iflangja{物理}{physical}};
  \draw[sg] (65,0)  rectangle (79,9)  node[midway, lab] {\iflangja{コード}{code}};
  \draw[sg] (65,9)  rectangle (79,14) node[midway, lab] {\iflangja{データ}{data}};
  \draw[sg] (65,14) rectangle (79,22) node[midway, lab] {\iflangja{ヒープ}{heap}};
  \draw[fill=white] (65,22) rectangle (79,30);
  \draw[sg] (65,30) rectangle (79,36) node[midway, lab] {\iflangja{スタック}{stack}};
  % physical memory
  \draw[oth] (101,-11) rectangle (115,36);
  \draw[sg] (101,-9) rectangle (115,-3) node[midway, lab] {\iflangja{スタック}{stack}};
  \draw[sg] (101,2)  rectangle (115,11) node[midway, lab] {\iflangja{コード}{code}};
  \draw[sg] (101,16) rectangle (115,24) node[midway, lab] {\iflangja{ヒープ}{heap}};
  \draw[sg] (101,28) rectangle (115,33) node[midway, lab] {\iflangja{データ}{data}};
  \draw[->] (79,4.5) -- (101,6.5);
  \draw[->] (79,11.5) -- (101,30.5);
  \draw[->] (79,18) -- (101,20);
  \draw[->] (79,33) -- (101,-6);
  \node[capt] at (90,-15) {\iflangja{(b) セグメンテーション}{(b) segmentation}};
\end{tikzpicture}

  \caption{(a) Paging maps fixed-size pages to page frames; page~2 is on
    disk, and the unused page has no mapping.  (b) Segmentation maps
    segments of different sizes to regions of physical memory of the same
    sizes.}
  \label{fig:l08-goals}
\end{figure}

\section{Hardware support}
\label{sec:l08-hardware}

Translating and validating every memory access in software would be far
too slow.\source{P3L2, memory management: hardware support; OSTEP ch.~15,
18--19.}  Memory management therefore relies on hardware in the CPU
package, which handles every access, while the operating system sets up the
data structures and handles the exceptions (\cref{fig:l08-mmu}).

\begin{definition}[Memory management unit]
  The \term{memory management unit}[MMU]\index{MMU|see{memory management unit}}
  (MMU) is the hardware component of the CPU package that translates the
  virtual addresses issued by the CPU into physical addresses and that
  raises a fault when a translation is not possible.
\end{definition}

The MMU reports a fault to the operating system in three situations: an
access to a virtual address that has not been allocated (an illegal
access); an access that violates the permissions, such as a write to
read-only memory; and an access to a page that is not present in physical
memory and must first be brought in from disk.  The MMU uses registers that
point to the page table of the running process, or, with segmentation,
that hold the base address, the limit and further information of the
segments.  Because walking a page table in memory is slow, the MMU also
contains a cache of valid translations, the TLB (\cref{sec:l08-tlb}).

The translation itself is performed entirely in hardware.  The operating
system creates and maintains the data structures that the hardware uses,
such as the page tables, but the hardware dictates their format and
therefore which memory management mechanisms are possible: for example,
the page sizes and the layout of a page table entry.

\begin{figure}[htbp]
  \centering
  % Hardware support for memory management: the MMU in the CPU package with
% its TLB and page table base register translates virtual addresses,
% walks the page tables in memory, and raises faults to the OS.
% Usage: % Hardware support for memory management: the MMU in the CPU package with
% its TLB and page table base register translates virtual addresses,
% walks the page tables in memory, and raises faults to the OS.
% Usage: % Hardware support for memory management: the MMU in the CPU package with
% its TLB and page table base register translates virtual addresses,
% walks the page tables in memory, and raises faults to the OS.
% Usage: \input{figures/l08-mmu}   (width ~116 mm)
\begin{tikzpicture}[x=1mm, y=1mm, font=\scriptsize\sffamily,
  >={Stealth[length=1.5mm]},
  box/.style={draw, fill=white, align=center, inner sep=1pt},
  hw/.style={draw, fill=ln-accent!14, align=center, inner sep=1pt},
  lab/.style={font=\scriptsize\sffamily, align=center},
  ttl/.style={font=\scriptsize\sffamily\bfseries, anchor=north west}]
  % CPU package
  \draw[ln-rule, dashed, rounded corners=3pt] (0,0) rectangle (64,36);
  \node[ttl, text=ln-muted] at (1,35.5) {\iflangja{CPU パッケージ}{CPU package}};
  \node[box, minimum width=14mm, minimum height=12mm] (cpu) at (10,17) {CPU};
  \draw[fill=ln-box] (30,3) rectangle (61,31);
  \node[ttl] at (30.5,30.5) {MMU};
  \node[hw, minimum width=27mm, minimum height=7mm] (tlb) at (45.5,19)
    {\iflangja{TLB（変換のキャッシュ）}{TLB (translation cache)}};
  \node[hw, minimum width=27mm, minimum height=7mm] (reg) at (45.5,9)
    {\iflangja{ページテーブル\\ベースレジスタ}{page table base\\register}};
  \draw[->] (cpu.east) -- (30,17) node[pos=0.5, above, lab] {\iflangja{仮想\\アドレス}{virtual\\address}};
  % physical memory
  \draw[fill=ln-box] (86,0) rectangle (116,36);
  \node[ttl] at (86.5,35.5) {\iflangja{物理メモリ}{physical memory}};
  \node[hw, minimum width=24mm, minimum height=7mm] (pt) at (101,24)
    {\iflangja{ページテーブル}{page tables}};
  \node[box, minimum width=24mm, minimum height=11mm] (pages) at (101,8)
    {\iflangja{プロセスの\\ページ}{pages of\\processes}};
  \draw[->] (61,8) -- (pages.west) node[pos=0.5, above, lab] {\iflangja{物理アドレス}{physical address}};
  \draw[<->, densely dashed] (61,24) -- (pt.west) node[pos=0.5, above, lab] {\iflangja{ページテーブル参照}{page table walk}};
  % faults
  \node[box, minimum width=31mm, minimum height=7mm, fill=ln-caution!8] (os) at (45.5,-12)
    {\iflangja{OS のフォールトハンドラ}{OS fault handler}};
  \draw[->, ln-caution] (45.5,3) -- (os.north) node[pos=0.45, right, lab, text=ln-caution]
    {\iflangja{フォールト}{fault}};
\end{tikzpicture}
   (width ~116 mm)
\begin{tikzpicture}[x=1mm, y=1mm, font=\scriptsize\sffamily,
  >={Stealth[length=1.5mm]},
  box/.style={draw, fill=white, align=center, inner sep=1pt},
  hw/.style={draw, fill=ln-accent!14, align=center, inner sep=1pt},
  lab/.style={font=\scriptsize\sffamily, align=center},
  ttl/.style={font=\scriptsize\sffamily\bfseries, anchor=north west}]
  % CPU package
  \draw[ln-rule, dashed, rounded corners=3pt] (0,0) rectangle (64,36);
  \node[ttl, text=ln-muted] at (1,35.5) {\iflangja{CPU パッケージ}{CPU package}};
  \node[box, minimum width=14mm, minimum height=12mm] (cpu) at (10,17) {CPU};
  \draw[fill=ln-box] (30,3) rectangle (61,31);
  \node[ttl] at (30.5,30.5) {MMU};
  \node[hw, minimum width=27mm, minimum height=7mm] (tlb) at (45.5,19)
    {\iflangja{TLB（変換のキャッシュ）}{TLB (translation cache)}};
  \node[hw, minimum width=27mm, minimum height=7mm] (reg) at (45.5,9)
    {\iflangja{ページテーブル\\ベースレジスタ}{page table base\\register}};
  \draw[->] (cpu.east) -- (30,17) node[pos=0.5, above, lab] {\iflangja{仮想\\アドレス}{virtual\\address}};
  % physical memory
  \draw[fill=ln-box] (86,0) rectangle (116,36);
  \node[ttl] at (86.5,35.5) {\iflangja{物理メモリ}{physical memory}};
  \node[hw, minimum width=24mm, minimum height=7mm] (pt) at (101,24)
    {\iflangja{ページテーブル}{page tables}};
  \node[box, minimum width=24mm, minimum height=11mm] (pages) at (101,8)
    {\iflangja{プロセスの\\ページ}{pages of\\processes}};
  \draw[->] (61,8) -- (pages.west) node[pos=0.5, above, lab] {\iflangja{物理アドレス}{physical address}};
  \draw[<->, densely dashed] (61,24) -- (pt.west) node[pos=0.5, above, lab] {\iflangja{ページテーブル参照}{page table walk}};
  % faults
  \node[box, minimum width=31mm, minimum height=7mm, fill=ln-caution!8] (os) at (45.5,-12)
    {\iflangja{OS のフォールトハンドラ}{OS fault handler}};
  \draw[->, ln-caution] (45.5,3) -- (os.north) node[pos=0.45, right, lab, text=ln-caution]
    {\iflangja{フォールト}{fault}};
\end{tikzpicture}
   (width ~116 mm)
\begin{tikzpicture}[x=1mm, y=1mm, font=\scriptsize\sffamily,
  >={Stealth[length=1.5mm]},
  box/.style={draw, fill=white, align=center, inner sep=1pt},
  hw/.style={draw, fill=ln-accent!14, align=center, inner sep=1pt},
  lab/.style={font=\scriptsize\sffamily, align=center},
  ttl/.style={font=\scriptsize\sffamily\bfseries, anchor=north west}]
  % CPU package
  \draw[ln-rule, dashed, rounded corners=3pt] (0,0) rectangle (64,36);
  \node[ttl, text=ln-muted] at (1,35.5) {\iflangja{CPU パッケージ}{CPU package}};
  \node[box, minimum width=14mm, minimum height=12mm] (cpu) at (10,17) {CPU};
  \draw[fill=ln-box] (30,3) rectangle (61,31);
  \node[ttl] at (30.5,30.5) {MMU};
  \node[hw, minimum width=27mm, minimum height=7mm] (tlb) at (45.5,19)
    {\iflangja{TLB（変換のキャッシュ）}{TLB (translation cache)}};
  \node[hw, minimum width=27mm, minimum height=7mm] (reg) at (45.5,9)
    {\iflangja{ページテーブル\\ベースレジスタ}{page table base\\register}};
  \draw[->] (cpu.east) -- (30,17) node[pos=0.5, above, lab] {\iflangja{仮想\\アドレス}{virtual\\address}};
  % physical memory
  \draw[fill=ln-box] (86,0) rectangle (116,36);
  \node[ttl] at (86.5,35.5) {\iflangja{物理メモリ}{physical memory}};
  \node[hw, minimum width=24mm, minimum height=7mm] (pt) at (101,24)
    {\iflangja{ページテーブル}{page tables}};
  \node[box, minimum width=24mm, minimum height=11mm] (pages) at (101,8)
    {\iflangja{プロセスの\\ページ}{pages of\\processes}};
  \draw[->] (61,8) -- (pages.west) node[pos=0.5, above, lab] {\iflangja{物理アドレス}{physical address}};
  \draw[<->, densely dashed] (61,24) -- (pt.west) node[pos=0.5, above, lab] {\iflangja{ページテーブル参照}{page table walk}};
  % faults
  \node[box, minimum width=31mm, minimum height=7mm, fill=ln-caution!8] (os) at (45.5,-12)
    {\iflangja{OS のフォールトハンドラ}{OS fault handler}};
  \draw[->, ln-caution] (45.5,3) -- (os.north) node[pos=0.45, right, lab, text=ln-caution]
    {\iflangja{フォールト}{fault}};
\end{tikzpicture}

  \caption{Hardware support for memory management.  The MMU translates
    virtual addresses with the help of its TLB, walks the page tables in
    physical memory when a translation is not cached, and raises a fault to
    the operating system when an access is not possible.}
  \label{fig:l08-mmu}
\end{figure}

\section{Page tables}
\label{sec:l08-page-tables}

A page table (Lesson~2) maps every page of a virtual address space to the
page frame that holds it.\source{P3L2, page tables; OSTEP ch.~18;
Silberschatz et al.\ ch.~9.}\index{page table}  Because pages and page
frames have the same size, the page table need not map every address
individually: the position of an address within its page is the same in
virtual and in physical memory, so only the page must be translated.

For pages of $2^{d}$ bytes, a virtual address is split into two parts.
\begin{description}
  \item[Virtual page number] The high-order bits form the
    \term{virtual page number}[仮想ページ番号]\index{VPN|see{virtual page number}}
    (VPN), which indexes the page table.
  \item[Offset] The $d$ low-order bits form the \term{offset}[オフセット],
    the position of the address within its page, which is not translated.
\end{description}
Each entry of the page table is a
\term{page table entry}[ページテーブルエントリ]\index{PTE|see{page table entry}}
(PTE).  Among other information (\cref{sec:l08-pte}), it holds the
\term{physical frame number}[物理フレーム番号]\index{PFN|see{physical frame number}}
(PFN) of the frame that contains the page.  The physical address combines
the PFN with the unchanged offset (\cref{fig:l08-translation}):
\begin{equation}
  \text{VPN} = \left\lfloor \frac{\mathit{VA}}{2^{d}} \right\rfloor ,
  \qquad
  \text{offset} = \mathit{VA} \bmod 2^{d} ,
  \qquad
  \mathit{PA} = \text{PFN} \cdot 2^{d} + \text{offset} .
  \label{eq:l08-translation}
\end{equation}

\begin{figure}[htbp]
  \centering
  % Address translation with a (flat) page table: the VPN indexes the page
% table found through the page table base register; the PFN of the entry
% is combined with the unchanged offset.
% Usage: % Address translation with a (flat) page table: the VPN indexes the page
% table found through the page table base register; the PFN of the entry
% is combined with the unchanged offset.
% Usage: % Address translation with a (flat) page table: the VPN indexes the page
% table found through the page table base register; the PFN of the entry
% is combined with the unchanged offset.
% Usage: \input{figures/l08-translation}   (width ~116 mm)
\begin{tikzpicture}[x=1mm, y=1mm, font=\scriptsize\sffamily,
  >={Stealth[length=1.5mm]},
  cell/.style={draw, minimum height=7mm, inner sep=0pt, anchor=south west},
  lab/.style={font=\scriptsize\sffamily, align=center},
  hd/.style={font=\scriptsize\sffamily\bfseries, anchor=south}]
  % virtual address
  \node[hd, anchor=south west] at (0,37.5) {\iflangja{仮想アドレス}{virtual address}};
  \node[cell, minimum width=18mm, fill=ln-accent!22] (vpn) at (0,30) {VPN};
  \node[cell, minimum width=16mm, fill=ln-box] (voff) at (18,30) {\iflangja{オフセット}{offset}};
  % page table base register
  \node[cell, minimum width=24mm, minimum height=6mm, fill=ln-source!14, align=center]
    (ptbr) at (46,36) {\iflangja{ページテーブル\\ベースレジスタ}{page table\\base register}};
  % page table
  \foreach \k in {0,...,5} {
    \draw[fill=white] (46,5*\k) rectangle ++(16,5);
    \draw[fill=white] (62,5*\k) rectangle ++(8,5);
  }
  \draw[fill=ln-accent!22] (46,15) rectangle ++(16,5) node[midway] {PFN};
  \draw[fill=ln-accent!10] (62,15) rectangle ++(8,5) node[midway, font=\iflangja{\tiny}{\scriptsize}\sffamily] {\iflangja{フラグ}{flags}};
  \node[lab, anchor=north] at (58,-0.5) {\iflangja{ページテーブル}{page table}};
  \draw[->] (58,36) -- (58,30);
  % VPN selects the entry
  \draw[->] (9,30) |- (46,17.5) node[pos=0.75, above, lab] {\iflangja{索引}{index}};
  % physical address
  \node[hd, anchor=north] at (99,13.5) {\iflangja{物理アドレス}{physical address}};
  \node[cell, minimum width=18mm, fill=ln-accent!22] (pfn) at (82,14) {PFN};
  \node[cell, minimum width=16mm, fill=ln-box] (poff) at (100,14) {\iflangja{オフセット}{offset}};
  \draw[->] (70,17.5) -- (82,17.5);
  % offset copied unchanged
  \draw[->, densely dashed] (26,37) -- (26,49) -| (108,21)
    node[pos=0.25, above, lab] {\iflangja{オフセットはそのまま}{offset copied unchanged}};
\end{tikzpicture}
   (width ~116 mm)
\begin{tikzpicture}[x=1mm, y=1mm, font=\scriptsize\sffamily,
  >={Stealth[length=1.5mm]},
  cell/.style={draw, minimum height=7mm, inner sep=0pt, anchor=south west},
  lab/.style={font=\scriptsize\sffamily, align=center},
  hd/.style={font=\scriptsize\sffamily\bfseries, anchor=south}]
  % virtual address
  \node[hd, anchor=south west] at (0,37.5) {\iflangja{仮想アドレス}{virtual address}};
  \node[cell, minimum width=18mm, fill=ln-accent!22] (vpn) at (0,30) {VPN};
  \node[cell, minimum width=16mm, fill=ln-box] (voff) at (18,30) {\iflangja{オフセット}{offset}};
  % page table base register
  \node[cell, minimum width=24mm, minimum height=6mm, fill=ln-source!14, align=center]
    (ptbr) at (46,36) {\iflangja{ページテーブル\\ベースレジスタ}{page table\\base register}};
  % page table
  \foreach \k in {0,...,5} {
    \draw[fill=white] (46,5*\k) rectangle ++(16,5);
    \draw[fill=white] (62,5*\k) rectangle ++(8,5);
  }
  \draw[fill=ln-accent!22] (46,15) rectangle ++(16,5) node[midway] {PFN};
  \draw[fill=ln-accent!10] (62,15) rectangle ++(8,5) node[midway, font=\iflangja{\tiny}{\scriptsize}\sffamily] {\iflangja{フラグ}{flags}};
  \node[lab, anchor=north] at (58,-0.5) {\iflangja{ページテーブル}{page table}};
  \draw[->] (58,36) -- (58,30);
  % VPN selects the entry
  \draw[->] (9,30) |- (46,17.5) node[pos=0.75, above, lab] {\iflangja{索引}{index}};
  % physical address
  \node[hd, anchor=north] at (99,13.5) {\iflangja{物理アドレス}{physical address}};
  \node[cell, minimum width=18mm, fill=ln-accent!22] (pfn) at (82,14) {PFN};
  \node[cell, minimum width=16mm, fill=ln-box] (poff) at (100,14) {\iflangja{オフセット}{offset}};
  \draw[->] (70,17.5) -- (82,17.5);
  % offset copied unchanged
  \draw[->, densely dashed] (26,37) -- (26,49) -| (108,21)
    node[pos=0.25, above, lab] {\iflangja{オフセットはそのまま}{offset copied unchanged}};
\end{tikzpicture}
   (width ~116 mm)
\begin{tikzpicture}[x=1mm, y=1mm, font=\scriptsize\sffamily,
  >={Stealth[length=1.5mm]},
  cell/.style={draw, minimum height=7mm, inner sep=0pt, anchor=south west},
  lab/.style={font=\scriptsize\sffamily, align=center},
  hd/.style={font=\scriptsize\sffamily\bfseries, anchor=south}]
  % virtual address
  \node[hd, anchor=south west] at (0,37.5) {\iflangja{仮想アドレス}{virtual address}};
  \node[cell, minimum width=18mm, fill=ln-accent!22] (vpn) at (0,30) {VPN};
  \node[cell, minimum width=16mm, fill=ln-box] (voff) at (18,30) {\iflangja{オフセット}{offset}};
  % page table base register
  \node[cell, minimum width=24mm, minimum height=6mm, fill=ln-source!14, align=center]
    (ptbr) at (46,36) {\iflangja{ページテーブル\\ベースレジスタ}{page table\\base register}};
  % page table
  \foreach \k in {0,...,5} {
    \draw[fill=white] (46,5*\k) rectangle ++(16,5);
    \draw[fill=white] (62,5*\k) rectangle ++(8,5);
  }
  \draw[fill=ln-accent!22] (46,15) rectangle ++(16,5) node[midway] {PFN};
  \draw[fill=ln-accent!10] (62,15) rectangle ++(8,5) node[midway, font=\iflangja{\tiny}{\scriptsize}\sffamily] {\iflangja{フラグ}{flags}};
  \node[lab, anchor=north] at (58,-0.5) {\iflangja{ページテーブル}{page table}};
  \draw[->] (58,36) -- (58,30);
  % VPN selects the entry
  \draw[->] (9,30) |- (46,17.5) node[pos=0.75, above, lab] {\iflangja{索引}{index}};
  % physical address
  \node[hd, anchor=north] at (99,13.5) {\iflangja{物理アドレス}{physical address}};
  \node[cell, minimum width=18mm, fill=ln-accent!22] (pfn) at (82,14) {PFN};
  \node[cell, minimum width=16mm, fill=ln-box] (poff) at (100,14) {\iflangja{オフセット}{offset}};
  \draw[->] (70,17.5) -- (82,17.5);
  % offset copied unchanged
  \draw[->, densely dashed] (26,37) -- (26,49) -| (108,21)
    node[pos=0.25, above, lab] {\iflangja{オフセットはそのまま}{offset copied unchanged}};
\end{tikzpicture}

  \caption{Address translation with a page table.  The VPN selects the
    entry of the page table of the running process; its PFN, combined with
    the unchanged offset, forms the physical address.}
  \label{fig:l08-translation}
\end{figure}

\begin{example}
  With 32-bit virtual addresses and \SI{4}{\kibi\byte} pages ($d = 12$), a
  VPN has 20 bits and an offset 12 bits, i.e., three hexadecimal digits.
  The virtual address \texttt{0x00035A7C} has the VPN \texttt{0x00035} and
  the offset \texttt{0xA7C}.  If entry \texttt{0x35} of the page table holds
  the PFN \texttt{0x001C2}, the physical address is \texttt{0x001C2A7C}.
  Every address from \texttt{0x00035000} to \texttt{0x00035FFF} is
  translated with the same entry.
\end{example}

\subsection{Allocation on first touch and reclaiming}

Physical memory is not allocated when a process reserves a range of virtual
addresses, for example for a large array, but only when the process
accesses the memory.  This policy is called
\term{allocation on first touch}[初回アクセス時の割り当て]: the first
access to a page faults, and the operating system then allocates a page
frame for the page and establishes the mapping in the page table.

\needspace{8\baselineskip}
\begin{example}
  The following code reserves \SI{256}{\mebi\byte} but receives page frames
  only as the loop touches the $2^{28}/2^{12} = 65\,536$ pages one after
  another:
  \begin{lstlisting}[language=C, numbers=none]
size_t n = 256UL << 20;          /* 256 MiB            */
char *buf = malloc(n);           /* no page frames yet */
for (size_t i = 0; i < n; i += 4096)
    buf[i] = 1;                  /* first touch faults */
\end{lstlisting}
\end{example}

Conversely, when a page has not been used for a long time, the operating
system may reclaim its frame: it moves the contents of the page to disk and
marks the mapping as invalid (\cref{sec:l08-pte}).  If the process accesses
the page again, the MMU raises a fault and traps into the operating system.
The operating system verifies that the access is allowed, brings the
contents back into memory, possibly into a different page frame, and
re-establishes the mapping; the process then continues without noticing
anything but the delay.

The operating system maintains one page table for every process.  The MMU
finds the page table of the running process through the
\term{page table base register}[ページテーブルベースレジスタ], a CPU
register that holds the address of that page table; on x86 it is the
register CR3.  On a context switch, the operating system loads the address
of the page table of the next process into this register.

\section{Page table entries}
\label{sec:l08-pte}

Besides the PFN,\source{P3L2, page table entry; OSTEP ch.~18, 21;
Silberschatz et al.\ ch.~9.} every PTE contains flags, which the MMU checks
on every access and which the operating system uses to manage memory.  The
most important flags are the following.
\begin{description}
  \item[Valid bit] The \term{valid bit}[有効ビット], also called the
    \emph{present bit}, indicates whether the mapping is valid, i.e.,
    whether the page is present in physical memory in the frame that the
    entry names.
  \item[Dirty bit] The hardware sets the \term{dirty bit}[ダーティビット]
    whenever the page is written.  Pages that cache the contents of a file,
    for example, need to be written back to the file only if their dirty
    bit is set.
  \item[Access bit] The hardware sets the \term{access bit}[アクセスビット]
    (referenced bit) whenever the page is accessed, for reading or for
    writing.
  \item[Protection bits] The \term{protection bits}[保護ビット] specify
    which operations are permitted on the page: read, write, execute.
\end{description}
The operating system reads the dirty and access bits and clears them when
it has acted on them.

On 32-bit x86 processors a PTE has 32 bits (\cref{fig:l08-pte}).  The upper
20 bits hold the PFN; among the lower bits are the present bit P, the dirty
bit D and the accessed bit A.  R/W is a protection bit: with 0 the page is
read-only, with 1 it may also be written.  U/S restricts the privilege
level: with 0 the page is accessible only in supervisor (kernel) mode, with
1 also in user mode.  PWT and PCD carry caching-related information
(write-through, cache disabled), and bits 9--11 are unused by the hardware
and available to the operating system.

\begin{figure}[htbp]
  \centering
  % Layout of a 32-bit x86 page table entry (4 KiB page, no PAE).
% Usage: % Layout of a 32-bit x86 page table entry (4 KiB page, no PAE).
% Usage: % Layout of a 32-bit x86 page table entry (4 KiB page, no PAE).
% Usage: \input{figures/l08-pte}   (width ~116 mm)
\begin{tikzpicture}[x=1mm, y=1mm, font=\scriptsize\sffamily,
  num/.style={font=\tiny\sffamily, text=ln-muted, anchor=south},
  grp/.style={font=\scriptsize\sffamily, anchor=north, align=center}]
  % field boundaries: PFN 0..34, unused 34..46, then 9 single bits of 7.8 mm
  \draw[fill=ln-accent!22] (0,0) rectangle (34,7) node[midway] {\iflangja{物理フレーム番号}{physical frame number}};
  \draw[fill=white] (34,0) rectangle (46,7) node[midway] {\iflangja{未使用}{unused}};
  \foreach \k/\n/\f in {0/G/white, 1/PAT/white, 2/D/ln-accent!22, 3/A/ln-accent!22,
                        4/PCD/ln-box, 5/PWT/ln-box, 6/{U/S}/ln-source!16,
                        7/{R/W}/ln-source!16, 8/P/ln-accent!22} {
    \draw[fill=\f] (46+7.8*\k,0) rectangle ++(7.8,7);
    \node at (49.9+7.8*\k,3.5) {\n};
  }
  % bit numbers
  \node[num] at (1.5,7.2) {31};
  \node[num] at (32.5,7.2) {12};
  \node[num] at (35.5,7.2) {11};
  \node[num] at (44.5,7.2) {9};
  \foreach \k/\b in {0/8, 1/7, 2/6, 3/5, 4/4, 5/3, 6/2, 7/1, 8/0} {
    \node[num] at (49.9+7.8*\k,7.2) {\b};
  }
  % groups
  \draw[ln-muted] (62.2,-1) -- (62.2,-2) -- (76.6,-2) -- (76.6,-1);
  \node[grp] at (69.4,-2.5) {\iflangja{ダーティ，\\アクセス}{dirty,\\accessed}};
  \draw[ln-muted] (77.8,-1) -- (77.8,-2) -- (92.2,-2) -- (92.2,-1);
  \node[grp] at (85,-2.5) {\iflangja{キャッシュ\\関連}{caching}};
  \draw[ln-muted] (93.4,-1) -- (93.4,-2) -- (107.8,-2) -- (107.8,-1);
  \node[grp] at (100.6,-2.5) {\iflangja{保護}{protection}};
  \draw[ln-muted] (112.3,-1) -- (112.3,-2);
  \node[grp] at (112.3,-2.5) {\iflangja{存在}{present}};
\end{tikzpicture}
   (width ~116 mm)
\begin{tikzpicture}[x=1mm, y=1mm, font=\scriptsize\sffamily,
  num/.style={font=\tiny\sffamily, text=ln-muted, anchor=south},
  grp/.style={font=\scriptsize\sffamily, anchor=north, align=center}]
  % field boundaries: PFN 0..34, unused 34..46, then 9 single bits of 7.8 mm
  \draw[fill=ln-accent!22] (0,0) rectangle (34,7) node[midway] {\iflangja{物理フレーム番号}{physical frame number}};
  \draw[fill=white] (34,0) rectangle (46,7) node[midway] {\iflangja{未使用}{unused}};
  \foreach \k/\n/\f in {0/G/white, 1/PAT/white, 2/D/ln-accent!22, 3/A/ln-accent!22,
                        4/PCD/ln-box, 5/PWT/ln-box, 6/{U/S}/ln-source!16,
                        7/{R/W}/ln-source!16, 8/P/ln-accent!22} {
    \draw[fill=\f] (46+7.8*\k,0) rectangle ++(7.8,7);
    \node at (49.9+7.8*\k,3.5) {\n};
  }
  % bit numbers
  \node[num] at (1.5,7.2) {31};
  \node[num] at (32.5,7.2) {12};
  \node[num] at (35.5,7.2) {11};
  \node[num] at (44.5,7.2) {9};
  \foreach \k/\b in {0/8, 1/7, 2/6, 3/5, 4/4, 5/3, 6/2, 7/1, 8/0} {
    \node[num] at (49.9+7.8*\k,7.2) {\b};
  }
  % groups
  \draw[ln-muted] (62.2,-1) -- (62.2,-2) -- (76.6,-2) -- (76.6,-1);
  \node[grp] at (69.4,-2.5) {\iflangja{ダーティ，\\アクセス}{dirty,\\accessed}};
  \draw[ln-muted] (77.8,-1) -- (77.8,-2) -- (92.2,-2) -- (92.2,-1);
  \node[grp] at (85,-2.5) {\iflangja{キャッシュ\\関連}{caching}};
  \draw[ln-muted] (93.4,-1) -- (93.4,-2) -- (107.8,-2) -- (107.8,-1);
  \node[grp] at (100.6,-2.5) {\iflangja{保護}{protection}};
  \draw[ln-muted] (112.3,-1) -- (112.3,-2);
  \node[grp] at (112.3,-2.5) {\iflangja{存在}{present}};
\end{tikzpicture}
   (width ~116 mm)
\begin{tikzpicture}[x=1mm, y=1mm, font=\scriptsize\sffamily,
  num/.style={font=\tiny\sffamily, text=ln-muted, anchor=south},
  grp/.style={font=\scriptsize\sffamily, anchor=north, align=center}]
  % field boundaries: PFN 0..34, unused 34..46, then 9 single bits of 7.8 mm
  \draw[fill=ln-accent!22] (0,0) rectangle (34,7) node[midway] {\iflangja{物理フレーム番号}{physical frame number}};
  \draw[fill=white] (34,0) rectangle (46,7) node[midway] {\iflangja{未使用}{unused}};
  \foreach \k/\n/\f in {0/G/white, 1/PAT/white, 2/D/ln-accent!22, 3/A/ln-accent!22,
                        4/PCD/ln-box, 5/PWT/ln-box, 6/{U/S}/ln-source!16,
                        7/{R/W}/ln-source!16, 8/P/ln-accent!22} {
    \draw[fill=\f] (46+7.8*\k,0) rectangle ++(7.8,7);
    \node at (49.9+7.8*\k,3.5) {\n};
  }
  % bit numbers
  \node[num] at (1.5,7.2) {31};
  \node[num] at (32.5,7.2) {12};
  \node[num] at (35.5,7.2) {11};
  \node[num] at (44.5,7.2) {9};
  \foreach \k/\b in {0/8, 1/7, 2/6, 3/5, 4/4, 5/3, 6/2, 7/1, 8/0} {
    \node[num] at (49.9+7.8*\k,7.2) {\b};
  }
  % groups
  \draw[ln-muted] (62.2,-1) -- (62.2,-2) -- (76.6,-2) -- (76.6,-1);
  \node[grp] at (69.4,-2.5) {\iflangja{ダーティ，\\アクセス}{dirty,\\accessed}};
  \draw[ln-muted] (77.8,-1) -- (77.8,-2) -- (92.2,-2) -- (92.2,-1);
  \node[grp] at (85,-2.5) {\iflangja{キャッシュ\\関連}{caching}};
  \draw[ln-muted] (93.4,-1) -- (93.4,-2) -- (107.8,-2) -- (107.8,-1);
  \node[grp] at (100.6,-2.5) {\iflangja{保護}{protection}};
  \draw[ln-muted] (112.3,-1) -- (112.3,-2);
  \node[grp] at (112.3,-2.5) {\iflangja{存在}{present}};
\end{tikzpicture}

  \caption{A 32-bit x86 page table entry.  The global bit G and the page
    attribute bit PAT are not discussed here.}
  \label{fig:l08-pte}
\end{figure}

\subsection{Page faults}
\label{sec:l08-page-fault}

The MMU uses a PTE both to translate an address and to check the
access.

\begin{definition}[Page fault]
  A \term{page fault}[ページフォールト] is the exception that the MMU
  raises when an access cannot be performed with the current page table
  entry: the entry is not valid, or the access violates its protection
  bits.
\end{definition}

On a page fault the CPU places an error code on the kernel stack and traps
into the kernel, which runs its page fault handler.  On x86 the error code
tells whether the page was not present or the access violated the
protection bits, whether the access was a read or a write, and whether it
came from user or kernel mode; the faulting virtual address is saved in
the register CR2.  From this information and from its own records of the
address space, the handler determines the action (\cref{tab:l08-faults}).

\begin{table}[htbp]
  \centering
  \caption{How the operating system handles an access, depending on the
    page table entry and on its own records.}
  \label{tab:l08-faults}
  \small
  \begin{tabular}{@{}lll@{}}
    \toprule
    Page table entry & Situation & Action \\
    \midrule
    valid          & access permitted              & none: the MMU translates \\
    not valid      & allocated, never accessed     & allocate a frame (first touch) \\
    not valid      & page is on disk               & bring it in (\cref{sec:l08-demand}) \\
    not valid      & address not allocated         & illegal access: \texttt{SIGSEGV} \\
    valid          & access violates protection    & \texttt{SIGSEGV}, or copy (\cref{sec:l08-cow}) \\
    \bottomrule
  \end{tabular}
\end{table}

A process that makes an illegal access receives the signal
\texttt{SIGSEGV} (Lesson~5), whose default action terminates it.  All other
faults are resolved by the operating system, after which the faulting
instruction is executed again.

\section{Page table size}
\label{sec:l08-pt-size}

A flat page table has one entry for every virtual page number, whether the
page is used or not.\source{P3L2, page table size; OSTEP ch.~20.}  For
$n$-bit virtual addresses, pages of $2^{d}$ bytes and page table entries of
$e$ bytes, the number of entries and the size of the table are
\begin{equation}
  N_{\mathrm{PTE}} = \frac{2^{n}}{2^{d}} = 2^{\,n-d} ,
  \qquad
  S_{\mathrm{PT}} = 2^{\,n-d} \cdot e .
  \label{eq:l08-pt-size}
\end{equation}
On a 32-bit architecture with \SI{4}{\kibi\byte} pages and 4-byte entries a
page table has $2^{20}$ entries and occupies \SI{4}{\mebi\byte}.  With
64-bit addresses and 8-byte entries it would have $2^{52}$ entries and
occupy $2^{55}$ bytes, \SI{32}{\pebi\byte}, for a single process.  Since
every process has its own page table, the cost is multiplied by the number
of processes.  Yet a process uses only a small part of its address space,
so almost all of these entries would describe pages that are never
used.\tip{Larger pages mean fewer entries: doubling the page size halves
$N_{\mathrm{PTE}}$ (\cref{sec:l08-page-size}).}

\begin{example}
  The x86-64 architecture uses 48-bit virtual addresses.  With
  \SI{4}{\kibi\byte} pages and 8-byte entries, a flat page table would need
  $2^{36}$ entries and $2^{39}$ bytes, \SI{512}{\gibi\byte}, per process.
\end{example}

\section{Multi-level page tables}
\label{sec:l08-multilevel}

Flat page tables are therefore impractical for large address
spaces.\source{P3L2, multi-level page tables; OSTEP ch.~20; Silberschatz et
al.\ ch.~9.}  Since most of an address space is unused, page tables are
instead organized hierarchically, so that no page table entries are needed
for the pages of large unused regions.

\begin{definition}[Multi-level page table]
  A \term{multi-level page table}[多段ページテーブル] is a tree of tables.
  The outer table, the \term{page directory}[ページディレクトリ], holds
  pointers to inner page tables, and the inner page tables hold the PFNs.
  An inner page table is allocated only for a region of the address space
  that contains valid pages.
\end{definition}

In a two-level table the virtual address is split into three parts
(\cref{fig:l08-multilevel}): $p_1$ indexes the page directory, $p_2$ the
inner page table that the directory entry points to, and the offset $d$
selects the address within the frame.  If $p_2$ has $k$ bits, one inner
page table covers $2^{k+d}$ bytes of the virtual address space.  A gap in
the address space of this size needs no inner page table at all; its
directory entry is simply marked invalid.  If the process later allocates
memory in that region, the operating system allocates an inner page table
for it and fills in the directory entry.

\begin{figure}[htbp]
  \centering
  % Two-level page table: p1 indexes the page directory, whose entry points
% to an inner page table; p2 indexes that table, whose entry gives the page
% frame; the offset d selects the byte within the frame.  Directory entries
% without an inner table are left white.
% Usage: % Two-level page table: p1 indexes the page directory, whose entry points
% to an inner page table; p2 indexes that table, whose entry gives the page
% frame; the offset d selects the byte within the frame.  Directory entries
% without an inner table are left white.
% Usage: % Two-level page table: p1 indexes the page directory, whose entry points
% to an inner page table; p2 indexes that table, whose entry gives the page
% frame; the offset d selects the byte within the frame.  Directory entries
% without an inner table are left white.
% Usage: \input{figures/l08-multilevel}   (width ~118 mm)
\begin{tikzpicture}[x=1mm, y=1mm, font=\scriptsize\sffamily,
  >={Stealth[length=1.5mm]},
  cell/.style={draw, minimum height=6mm, inner sep=0pt, anchor=south west},
  lab/.style={font=\scriptsize\sffamily, align=center},
  idx/.style={font=\scriptsize\sffamily, text=ln-muted},
  hd/.style={font=\scriptsize\sffamily\bfseries, anchor=south, align=center}]
  % page table base register
  \node[cell, minimum width=22mm, minimum height=7mm, fill=ln-source!14, align=center]
    at (0,38) {\iflangja{ページテーブル\\ベースレジスタ}{page table\\base register}};
  \draw[->] (22,41.5) -| (39,36);
  % page directory: rows top to bottom, centres 33, 27, 21, 15, 9
  \node[hd] at (39,44) {\iflangja{ページ\\ディレクトリ}{page\\directory}};
  \foreach \r/\f in {0/ln-accent!10, 1/ln-accent!30, 2/ln-accent!10, 3/white, 4/white} {
    \draw[fill=\f] (30,30-6*\r) rectangle ++(18,6);
  }
  \node[lab] at (39,27) {\iflangja{内部表へ}{$\to$ inner}};
  % inner page table
  \node[hd] at (73,44) {\iflangja{内部\\ページテーブル}{inner\\page table}};
  \foreach \r/\f in {0/ln-accent!10, 1/white, 2/ln-accent!10, 3/ln-accent!30, 4/white} {
    \draw[fill=\f] (64,30-6*\r) rectangle ++(18,6);
  }
  \node[lab] at (73,15) {PFN};
  % page frame
  \node[hd] at (107,44) {\iflangja{ページ\\フレーム}{page\\frame}};
  \draw[fill=ln-box] (98,6) rectangle (116,36);
  \draw[fill=ln-accent!30] (98,17) rectangle (116,20);
  % pointers along the top
  \draw[->] (48,27) -- (56,27) |- (73,40.5) -- (73,36);
  \draw[->] (82,15) -- (90,15) |- (107,40.5) -- (107,36);
  % virtual address
  \node[cell, minimum width=16mm, fill=ln-accent!22] (p1) at (31,-16) {$p_1$};
  \node[cell, minimum width=16mm, fill=ln-accent!22] (p2) at (47,-16) {$p_2$};
  \node[cell, minimum width=16mm, fill=ln-box] (d) at (63,-16) {$d$};
  \node[lab, anchor=east] at (30.5,-13) {\iflangja{仮想アドレス}{virtual\\address}};
  % indices from below
  \draw[->] (39,-10) -- (39,6);
  \draw[->] (55,-10) |- (73,-2) -- (73,6);
  \draw[->] (71,-10) |- (107,-6) -- (107,17);
  \node[idx, anchor=west] at (39.5,1) {\iflangja{索引}{index}};
  \node[idx, anchor=west] at (73.5,2) {\iflangja{索引}{index}};
  \node[idx, anchor=west] at (107.5,2) {\iflangja{オフセット}{offset}};
\end{tikzpicture}
   (width ~118 mm)
\begin{tikzpicture}[x=1mm, y=1mm, font=\scriptsize\sffamily,
  >={Stealth[length=1.5mm]},
  cell/.style={draw, minimum height=6mm, inner sep=0pt, anchor=south west},
  lab/.style={font=\scriptsize\sffamily, align=center},
  idx/.style={font=\scriptsize\sffamily, text=ln-muted},
  hd/.style={font=\scriptsize\sffamily\bfseries, anchor=south, align=center}]
  % page table base register
  \node[cell, minimum width=22mm, minimum height=7mm, fill=ln-source!14, align=center]
    at (0,38) {\iflangja{ページテーブル\\ベースレジスタ}{page table\\base register}};
  \draw[->] (22,41.5) -| (39,36);
  % page directory: rows top to bottom, centres 33, 27, 21, 15, 9
  \node[hd] at (39,44) {\iflangja{ページ\\ディレクトリ}{page\\directory}};
  \foreach \r/\f in {0/ln-accent!10, 1/ln-accent!30, 2/ln-accent!10, 3/white, 4/white} {
    \draw[fill=\f] (30,30-6*\r) rectangle ++(18,6);
  }
  \node[lab] at (39,27) {\iflangja{内部表へ}{$\to$ inner}};
  % inner page table
  \node[hd] at (73,44) {\iflangja{内部\\ページテーブル}{inner\\page table}};
  \foreach \r/\f in {0/ln-accent!10, 1/white, 2/ln-accent!10, 3/ln-accent!30, 4/white} {
    \draw[fill=\f] (64,30-6*\r) rectangle ++(18,6);
  }
  \node[lab] at (73,15) {PFN};
  % page frame
  \node[hd] at (107,44) {\iflangja{ページ\\フレーム}{page\\frame}};
  \draw[fill=ln-box] (98,6) rectangle (116,36);
  \draw[fill=ln-accent!30] (98,17) rectangle (116,20);
  % pointers along the top
  \draw[->] (48,27) -- (56,27) |- (73,40.5) -- (73,36);
  \draw[->] (82,15) -- (90,15) |- (107,40.5) -- (107,36);
  % virtual address
  \node[cell, minimum width=16mm, fill=ln-accent!22] (p1) at (31,-16) {$p_1$};
  \node[cell, minimum width=16mm, fill=ln-accent!22] (p2) at (47,-16) {$p_2$};
  \node[cell, minimum width=16mm, fill=ln-box] (d) at (63,-16) {$d$};
  \node[lab, anchor=east] at (30.5,-13) {\iflangja{仮想アドレス}{virtual\\address}};
  % indices from below
  \draw[->] (39,-10) -- (39,6);
  \draw[->] (55,-10) |- (73,-2) -- (73,6);
  \draw[->] (71,-10) |- (107,-6) -- (107,17);
  \node[idx, anchor=west] at (39.5,1) {\iflangja{索引}{index}};
  \node[idx, anchor=west] at (73.5,2) {\iflangja{索引}{index}};
  \node[idx, anchor=west] at (107.5,2) {\iflangja{オフセット}{offset}};
\end{tikzpicture}
   (width ~118 mm)
\begin{tikzpicture}[x=1mm, y=1mm, font=\scriptsize\sffamily,
  >={Stealth[length=1.5mm]},
  cell/.style={draw, minimum height=6mm, inner sep=0pt, anchor=south west},
  lab/.style={font=\scriptsize\sffamily, align=center},
  idx/.style={font=\scriptsize\sffamily, text=ln-muted},
  hd/.style={font=\scriptsize\sffamily\bfseries, anchor=south, align=center}]
  % page table base register
  \node[cell, minimum width=22mm, minimum height=7mm, fill=ln-source!14, align=center]
    at (0,38) {\iflangja{ページテーブル\\ベースレジスタ}{page table\\base register}};
  \draw[->] (22,41.5) -| (39,36);
  % page directory: rows top to bottom, centres 33, 27, 21, 15, 9
  \node[hd] at (39,44) {\iflangja{ページ\\ディレクトリ}{page\\directory}};
  \foreach \r/\f in {0/ln-accent!10, 1/ln-accent!30, 2/ln-accent!10, 3/white, 4/white} {
    \draw[fill=\f] (30,30-6*\r) rectangle ++(18,6);
  }
  \node[lab] at (39,27) {\iflangja{内部表へ}{$\to$ inner}};
  % inner page table
  \node[hd] at (73,44) {\iflangja{内部\\ページテーブル}{inner\\page table}};
  \foreach \r/\f in {0/ln-accent!10, 1/white, 2/ln-accent!10, 3/ln-accent!30, 4/white} {
    \draw[fill=\f] (64,30-6*\r) rectangle ++(18,6);
  }
  \node[lab] at (73,15) {PFN};
  % page frame
  \node[hd] at (107,44) {\iflangja{ページ\\フレーム}{page\\frame}};
  \draw[fill=ln-box] (98,6) rectangle (116,36);
  \draw[fill=ln-accent!30] (98,17) rectangle (116,20);
  % pointers along the top
  \draw[->] (48,27) -- (56,27) |- (73,40.5) -- (73,36);
  \draw[->] (82,15) -- (90,15) |- (107,40.5) -- (107,36);
  % virtual address
  \node[cell, minimum width=16mm, fill=ln-accent!22] (p1) at (31,-16) {$p_1$};
  \node[cell, minimum width=16mm, fill=ln-accent!22] (p2) at (47,-16) {$p_2$};
  \node[cell, minimum width=16mm, fill=ln-box] (d) at (63,-16) {$d$};
  \node[lab, anchor=east] at (30.5,-13) {\iflangja{仮想アドレス}{virtual\\address}};
  % indices from below
  \draw[->] (39,-10) -- (39,6);
  \draw[->] (55,-10) |- (73,-2) -- (73,6);
  \draw[->] (71,-10) |- (107,-6) -- (107,17);
  \node[idx, anchor=west] at (39.5,1) {\iflangja{索引}{index}};
  \node[idx, anchor=west] at (73.5,2) {\iflangja{索引}{index}};
  \node[idx, anchor=west] at (107.5,2) {\iflangja{オフセット}{offset}};
\end{tikzpicture}

  \caption{A two-level page table.  White directory entries have no inner
    page table; the translation takes one access to the page directory and
    one to the inner page table.}
  \label{fig:l08-multilevel}
\end{figure}

\begin{example}
  \label{ex:l08-multilevel}
  Consider 20-bit virtual addresses and \SI{1}{\kibi\byte} pages ($d =
  10$).  A flat page table has $2^{10} = 1024$ entries.  A two-level table
  with $p_1$ of 4 bits and $p_2$ of 6 bits has a page directory of 16
  entries, and each inner table of 64 entries covers \SI{64}{\kibi\byte}.  A
  process uses the lowest \SI{100}{\kibi\byte} (text, data and heap, pages
  0--99) and the highest \SI{20}{\kibi\byte} (stack, pages 1004--1023).  It
  needs inner tables for pages 0--63, 64--127 and 960--1023, in total $16 +
  3 \cdot 64 = 208$ entries, about a fifth of the flat table.
\end{example}

The hierarchy can be extended with further levels.  A third level adds a
table of pointers to page directories, and a fourth level a table of
pointers to those.\margin{x86-64 uses four levels of 9 bits each above a
12-bit offset: $4\cdot 9 + 12 = 48$ address bits.}  Additional levels matter
most on 64-bit architectures, whose address spaces are both larger and
sparser: the larger the region covered by a single entry high in the tree,
the larger the gaps that need no tables below it.

The price is translation latency.  Every level adds one memory access to
the translation of an address, since the entries of all levels are read
from memory in turn.  More levels save more memory for the page tables but
make each translation slower.

\section{Speeding up translation: the TLB}
\label{sec:l08-tlb}

With a single-level page table every memory reference requires two memory
accesses, one to the page table entry and one to the data; with four levels
it requires five.\source{P3L2, speed of translation: TLB; OSTEP ch.~19;
Silberschatz et al.\ ch.~9.}  The MMU avoids most of these accesses with a
cache.

\begin{definition}[Translation lookaside buffer]
  A
  \term{translation lookaside buffer}[TLB]\index{TLB|see{translation lookaside buffer}}
  (TLB) is a cache of recently used address translations integrated into
  the MMU.  On a TLB hit the address is translated without accessing the
  page table; on a TLB miss the page table is walked, and the translation
  is then stored in the TLB.
\end{definition}

A TLB entry also holds the protection and validity bits of the translation,
so that the access can be checked without the page table.  A small number
of entries suffices for a high hit rate because memory accesses exhibit
\emph{locality}: a page that has just been accessed is likely to be accessed
again soon (temporal locality), and addresses close to it, which usually lie
in the same page, are likely to be accessed next (spatial locality).  An
Intel Core i7 processor of the Sandy Bridge generation, for example, has per
core a data TLB of 64 entries, an instruction TLB of 128 entries, and a
second-level TLB of 512 entries shared by both kinds of translations (all
for \SI{4}{\kibi\byte} pages).\margin{A context switch makes TLB entries stale; they are flushed
unless tagged with their address space (Lesson~12).}

With TLB lookup time $t_{\mathrm{TLB}}$, memory access time $t_{\mathrm{mem}}$,
hit rate $h$ and an $L$-level page table, the effective time of a memory
access is
\begin{equation}
  t_{\mathrm{eff}} = h\,(t_{\mathrm{TLB}} + t_{\mathrm{mem}})
    + (1-h)\,\bigl(t_{\mathrm{TLB}} + (L+1)\,t_{\mathrm{mem}}\bigr) .
  \label{eq:l08-eat}
\end{equation}

\begin{example}
  Let $t_{\mathrm{mem}} = \SI{80}{\nano\second}$,
  $t_{\mathrm{TLB}} = \SI{1}{\nano\second}$ and $L = 4$.  Without a TLB,
  every access costs $5 \cdot 80 = \SI{400}{\nano\second}$.  With a hit rate
  of \SI{99}{\percent}, \cref{eq:l08-eat} gives $0.99 \cdot 81 + 0.01 \cdot
  401 \approx \SI{84.2}{\nano\second}$, only \SI{5}{\percent} more than an
  untranslated access.  With a hit rate of \SI{90}{\percent} it gives $0.9
  \cdot 81 + 0.1 \cdot 401 = \SI{113}{\nano\second}$.
\end{example}

\begin{keypoint}
  Flat page tables are proportional to the virtual address space and far
  too large.  Multi-level page tables allocate tables only for the regions
  in use, at the price of one memory access per level; the TLB, which
  exploits locality, avoids those accesses for almost every translation.
\end{keypoint}

\section{Inverted page tables}
\label{sec:l08-inverted}

A different way to keep page tables small is to organize them by physical
rather than by virtual memory.\source{P3L2, inverted page tables; OSTEP
ch.~20; Silberschatz et al.\ ch.~9.}  Physical memory is much smaller than
the virtual address spaces of all processes together, and every page frame
holds at most one page at a time.

\begin{definition}[Inverted page table]
  An \term{inverted page table}[逆ページテーブル] has one entry for every
  page frame of physical memory.  The entry of frame $i$ records the
  process identifier and the virtual page number of the page that frame
  $i$ holds.
\end{definition}

The system needs a single inverted page table, whose size is proportional
to physical memory and independent of the number of processes and of the
size of their address spaces.  A virtual address consists of the process
identifier (pid), the page number $p$ and the offset $d$.  To translate it,
the table is searched for the entry that matches $(\text{pid}, p)$; the
index $i$ of that entry is the frame number, and the physical address is
formed from $i$ and $d$ (\cref{fig:l08-inverted}\,(a)).

\begin{figure}[htbp]
  \centering
  % (a) Inverted page table: one entry per page frame; the entry matching
% (pid, p) is searched, and its index i is the frame number.
% (b) Hashed page table: a hash of (pid, p) selects a bucket whose chain
% holds (pid, p, frame) entries.
% Usage: % (a) Inverted page table: one entry per page frame; the entry matching
% (pid, p) is searched, and its index i is the frame number.
% (b) Hashed page table: a hash of (pid, p) selects a bucket whose chain
% holds (pid, p, frame) entries.
% Usage: % (a) Inverted page table: one entry per page frame; the entry matching
% (pid, p) is searched, and its index i is the frame number.
% (b) Hashed page table: a hash of (pid, p) selects a bucket whose chain
% holds (pid, p, frame) entries.
% Usage: \input{figures/l08-inverted}   (width ~118 mm)
\begin{tikzpicture}[x=1mm, y=1mm, font=\scriptsize\sffamily,
  >={Stealth[length=1.5mm]},
  cell/.style={draw, minimum height=6mm, inner sep=0pt, anchor=south west},
  lab/.style={font=\scriptsize\sffamily, align=center},
  idx/.style={font=\scriptsize\sffamily, text=ln-muted},
  hd/.style={font=\scriptsize\sffamily\bfseries, anchor=south, align=center},
  capt/.style={font=\footnotesize\sffamily, anchor=west}]
  % ================= (a) inverted page table
  \node[capt] at (0,42) {\iflangja{(a) 逆ページテーブル}{(a) inverted page table}};
  \node[lab, anchor=south west] at (0,27) {\iflangja{仮想アドレス}{virtual address}};
  \node[cell, minimum width=10mm, fill=ln-source!16] at (0,20) {pid};
  \node[cell, minimum width=10mm, fill=ln-accent!22] at (10,20) {$p$};
  \node[cell, minimum width=10mm, fill=ln-box] at (20,20) {$d$};
  \node[hd] at (58,31) {\iflangja{フレームごとに 1 エントリ}{one entry per frame}};
  \foreach \r/\a/\b in {0/7/3, 1/2/12, 2/5/9, 4/7/8} {
    \draw[fill=white] (48,24-6*\r) rectangle ++(10,6) node[midway] {\a};
    \draw[fill=white] (58,24-6*\r) rectangle ++(10,6) node[midway] {\b};
  }
  \draw[fill=ln-accent!22] (48,6) rectangle ++(10,6) node[midway] {pid};
  \draw[fill=ln-accent!22] (58,6) rectangle ++(10,6) node[midway] {$p$};
  \node[idx, anchor=east] at (47.5,9) {$i$};
  \node[idx] at (58,-1.5) {\iflangja{（索引 $i$ ＝ フレーム番号）}{(index $i$ = frame number)}};
  \draw (5,20) -- (5,14) -- (15,14);
  \draw[->] (15,20) |- (44,9) node[pos=0.75, above, lab] {\iflangja{探索}{search}};
  % physical address
  \node[cell, minimum width=10mm, fill=ln-accent!22] at (86,6) {$i$};
  \node[cell, minimum width=10mm, fill=ln-box] at (96,6) {$d$};
  \node[lab, anchor=north west] at (86,5.5) {\iflangja{物理アドレス}{physical address}};
  \draw[->] (68,9) -- (86,9);
  \draw[->, densely dashed] (25,26) -- (25,37) -| (101,12);
  % ================= (b) hashed page table
  \node[capt] at (0,-8) {\iflangja{(b) ハッシュページテーブル}{(b) hashed page table}};
  \node[cell, minimum width=10mm, fill=ln-source!16] at (0,-24) {pid};
  \node[cell, minimum width=10mm, fill=ln-accent!22] at (10,-24) {$p$};
  \node[cell, minimum width=10mm, fill=ln-box] at (20,-24) {$d$};
  \node[draw, circle, inner sep=1pt, fill=white, font=\scriptsize\sffamily] (h) at (34,-21) {$h$};
  \draw (5,-18) -- (5,-14) -- (34,-14);
  \draw (15,-18) -- (15,-14);
  \draw[->] (34,-14) -- (h.north);
  % hash table buckets
  \foreach \r in {0,...,4} { \draw[fill=white] (44,-12-5*\r) rectangle ++(8,5); }
  \draw[fill=ln-accent!22] (44,-27) rectangle ++(8,5);
  \node[idx] at (48,-34.5) {\iflangja{バケット}{buckets}};
  \draw[->] (h.east) -- (44,-24.5);
  % chain entries: pid | p | frame
  \foreach \x/\a/\b/\c in {60/4/31/17, 88/2/31/6} {
    \draw[fill=white] (\x,-27) rectangle ++(6,5) node[midway] {\a};
    \draw[fill=white] (\x+6,-27) rectangle ++(6,5) node[midway] {\b};
    \draw[fill=ln-accent!10] (\x+12,-27) rectangle ++(6,5) node[midway] {\c};
    \draw[fill=ln-box] (\x+18,-27) rectangle ++(3,5);
  }
  \draw[->] (48,-24.5) -- (60,-24.5);
  \fill (79.5,-24.5) circle (0.5);
  \draw[->] (79.5,-24.5) -- (88,-24.5);
  \node[idx, anchor=north] at (69,-27.5) {pid\quad $p$\quad \iflangja{枠}{frame}};
  \node[idx, anchor=north] at (97,-27.5) {pid\quad $p$\quad \iflangja{枠}{frame}};
  \node[lab, text width=40mm, anchor=north west] at (60,-33)
    {\iflangja{連鎖を (pid, $p$) で照合}{chain searched for (pid, $p$)}};
\end{tikzpicture}
   (width ~118 mm)
\begin{tikzpicture}[x=1mm, y=1mm, font=\scriptsize\sffamily,
  >={Stealth[length=1.5mm]},
  cell/.style={draw, minimum height=6mm, inner sep=0pt, anchor=south west},
  lab/.style={font=\scriptsize\sffamily, align=center},
  idx/.style={font=\scriptsize\sffamily, text=ln-muted},
  hd/.style={font=\scriptsize\sffamily\bfseries, anchor=south, align=center},
  capt/.style={font=\footnotesize\sffamily, anchor=west}]
  % ================= (a) inverted page table
  \node[capt] at (0,42) {\iflangja{(a) 逆ページテーブル}{(a) inverted page table}};
  \node[lab, anchor=south west] at (0,27) {\iflangja{仮想アドレス}{virtual address}};
  \node[cell, minimum width=10mm, fill=ln-source!16] at (0,20) {pid};
  \node[cell, minimum width=10mm, fill=ln-accent!22] at (10,20) {$p$};
  \node[cell, minimum width=10mm, fill=ln-box] at (20,20) {$d$};
  \node[hd] at (58,31) {\iflangja{フレームごとに 1 エントリ}{one entry per frame}};
  \foreach \r/\a/\b in {0/7/3, 1/2/12, 2/5/9, 4/7/8} {
    \draw[fill=white] (48,24-6*\r) rectangle ++(10,6) node[midway] {\a};
    \draw[fill=white] (58,24-6*\r) rectangle ++(10,6) node[midway] {\b};
  }
  \draw[fill=ln-accent!22] (48,6) rectangle ++(10,6) node[midway] {pid};
  \draw[fill=ln-accent!22] (58,6) rectangle ++(10,6) node[midway] {$p$};
  \node[idx, anchor=east] at (47.5,9) {$i$};
  \node[idx] at (58,-1.5) {\iflangja{（索引 $i$ ＝ フレーム番号）}{(index $i$ = frame number)}};
  \draw (5,20) -- (5,14) -- (15,14);
  \draw[->] (15,20) |- (44,9) node[pos=0.75, above, lab] {\iflangja{探索}{search}};
  % physical address
  \node[cell, minimum width=10mm, fill=ln-accent!22] at (86,6) {$i$};
  \node[cell, minimum width=10mm, fill=ln-box] at (96,6) {$d$};
  \node[lab, anchor=north west] at (86,5.5) {\iflangja{物理アドレス}{physical address}};
  \draw[->] (68,9) -- (86,9);
  \draw[->, densely dashed] (25,26) -- (25,37) -| (101,12);
  % ================= (b) hashed page table
  \node[capt] at (0,-8) {\iflangja{(b) ハッシュページテーブル}{(b) hashed page table}};
  \node[cell, minimum width=10mm, fill=ln-source!16] at (0,-24) {pid};
  \node[cell, minimum width=10mm, fill=ln-accent!22] at (10,-24) {$p$};
  \node[cell, minimum width=10mm, fill=ln-box] at (20,-24) {$d$};
  \node[draw, circle, inner sep=1pt, fill=white, font=\scriptsize\sffamily] (h) at (34,-21) {$h$};
  \draw (5,-18) -- (5,-14) -- (34,-14);
  \draw (15,-18) -- (15,-14);
  \draw[->] (34,-14) -- (h.north);
  % hash table buckets
  \foreach \r in {0,...,4} { \draw[fill=white] (44,-12-5*\r) rectangle ++(8,5); }
  \draw[fill=ln-accent!22] (44,-27) rectangle ++(8,5);
  \node[idx] at (48,-34.5) {\iflangja{バケット}{buckets}};
  \draw[->] (h.east) -- (44,-24.5);
  % chain entries: pid | p | frame
  \foreach \x/\a/\b/\c in {60/4/31/17, 88/2/31/6} {
    \draw[fill=white] (\x,-27) rectangle ++(6,5) node[midway] {\a};
    \draw[fill=white] (\x+6,-27) rectangle ++(6,5) node[midway] {\b};
    \draw[fill=ln-accent!10] (\x+12,-27) rectangle ++(6,5) node[midway] {\c};
    \draw[fill=ln-box] (\x+18,-27) rectangle ++(3,5);
  }
  \draw[->] (48,-24.5) -- (60,-24.5);
  \fill (79.5,-24.5) circle (0.5);
  \draw[->] (79.5,-24.5) -- (88,-24.5);
  \node[idx, anchor=north] at (69,-27.5) {pid\quad $p$\quad \iflangja{枠}{frame}};
  \node[idx, anchor=north] at (97,-27.5) {pid\quad $p$\quad \iflangja{枠}{frame}};
  \node[lab, text width=40mm, anchor=north west] at (60,-33)
    {\iflangja{連鎖を (pid, $p$) で照合}{chain searched for (pid, $p$)}};
\end{tikzpicture}
   (width ~118 mm)
\begin{tikzpicture}[x=1mm, y=1mm, font=\scriptsize\sffamily,
  >={Stealth[length=1.5mm]},
  cell/.style={draw, minimum height=6mm, inner sep=0pt, anchor=south west},
  lab/.style={font=\scriptsize\sffamily, align=center},
  idx/.style={font=\scriptsize\sffamily, text=ln-muted},
  hd/.style={font=\scriptsize\sffamily\bfseries, anchor=south, align=center},
  capt/.style={font=\footnotesize\sffamily, anchor=west}]
  % ================= (a) inverted page table
  \node[capt] at (0,42) {\iflangja{(a) 逆ページテーブル}{(a) inverted page table}};
  \node[lab, anchor=south west] at (0,27) {\iflangja{仮想アドレス}{virtual address}};
  \node[cell, minimum width=10mm, fill=ln-source!16] at (0,20) {pid};
  \node[cell, minimum width=10mm, fill=ln-accent!22] at (10,20) {$p$};
  \node[cell, minimum width=10mm, fill=ln-box] at (20,20) {$d$};
  \node[hd] at (58,31) {\iflangja{フレームごとに 1 エントリ}{one entry per frame}};
  \foreach \r/\a/\b in {0/7/3, 1/2/12, 2/5/9, 4/7/8} {
    \draw[fill=white] (48,24-6*\r) rectangle ++(10,6) node[midway] {\a};
    \draw[fill=white] (58,24-6*\r) rectangle ++(10,6) node[midway] {\b};
  }
  \draw[fill=ln-accent!22] (48,6) rectangle ++(10,6) node[midway] {pid};
  \draw[fill=ln-accent!22] (58,6) rectangle ++(10,6) node[midway] {$p$};
  \node[idx, anchor=east] at (47.5,9) {$i$};
  \node[idx] at (58,-1.5) {\iflangja{（索引 $i$ ＝ フレーム番号）}{(index $i$ = frame number)}};
  \draw (5,20) -- (5,14) -- (15,14);
  \draw[->] (15,20) |- (44,9) node[pos=0.75, above, lab] {\iflangja{探索}{search}};
  % physical address
  \node[cell, minimum width=10mm, fill=ln-accent!22] at (86,6) {$i$};
  \node[cell, minimum width=10mm, fill=ln-box] at (96,6) {$d$};
  \node[lab, anchor=north west] at (86,5.5) {\iflangja{物理アドレス}{physical address}};
  \draw[->] (68,9) -- (86,9);
  \draw[->, densely dashed] (25,26) -- (25,37) -| (101,12);
  % ================= (b) hashed page table
  \node[capt] at (0,-8) {\iflangja{(b) ハッシュページテーブル}{(b) hashed page table}};
  \node[cell, minimum width=10mm, fill=ln-source!16] at (0,-24) {pid};
  \node[cell, minimum width=10mm, fill=ln-accent!22] at (10,-24) {$p$};
  \node[cell, minimum width=10mm, fill=ln-box] at (20,-24) {$d$};
  \node[draw, circle, inner sep=1pt, fill=white, font=\scriptsize\sffamily] (h) at (34,-21) {$h$};
  \draw (5,-18) -- (5,-14) -- (34,-14);
  \draw (15,-18) -- (15,-14);
  \draw[->] (34,-14) -- (h.north);
  % hash table buckets
  \foreach \r in {0,...,4} { \draw[fill=white] (44,-12-5*\r) rectangle ++(8,5); }
  \draw[fill=ln-accent!22] (44,-27) rectangle ++(8,5);
  \node[idx] at (48,-34.5) {\iflangja{バケット}{buckets}};
  \draw[->] (h.east) -- (44,-24.5);
  % chain entries: pid | p | frame
  \foreach \x/\a/\b/\c in {60/4/31/17, 88/2/31/6} {
    \draw[fill=white] (\x,-27) rectangle ++(6,5) node[midway] {\a};
    \draw[fill=white] (\x+6,-27) rectangle ++(6,5) node[midway] {\b};
    \draw[fill=ln-accent!10] (\x+12,-27) rectangle ++(6,5) node[midway] {\c};
    \draw[fill=ln-box] (\x+18,-27) rectangle ++(3,5);
  }
  \draw[->] (48,-24.5) -- (60,-24.5);
  \fill (79.5,-24.5) circle (0.5);
  \draw[->] (79.5,-24.5) -- (88,-24.5);
  \node[idx, anchor=north] at (69,-27.5) {pid\quad $p$\quad \iflangja{枠}{frame}};
  \node[idx, anchor=north] at (97,-27.5) {pid\quad $p$\quad \iflangja{枠}{frame}};
  \node[lab, text width=40mm, anchor=north west] at (60,-33)
    {\iflangja{連鎖を (pid, $p$) で照合}{chain searched for (pid, $p$)}};
\end{tikzpicture}

  \caption{(a) An inverted page table is searched for $(\text{pid}, p)$;
    the index of the matching entry is the frame number.  (b) A hashed page
    table searches only the chain of the bucket selected by the hash
    function $h$.}
  \label{fig:l08-inverted}
\end{figure}

A linear search of the table is slow.  Most translations are served by the
TLB, but the search on a miss must still be fast, and it is accelerated
with hashing.  In a \term{hashed page table}[ハッシュページテーブル] a hash
function applied to the page number, together with the process identifier,
selects a bucket, which heads a linked list of the entries whose keys hash
to that bucket (\cref{fig:l08-inverted}\,(b)).  Each element holds the
process identifier, the page number, the frame number and a pointer to the
next element, so a translation searches only one short list.

\begin{example}
  A machine with \SI{16}{\gibi\byte} of physical memory and
  \SI{4}{\kibi\byte} frames has $2^{34}/2^{12} = 2^{22}$ frames.  With
  entries of 16 bytes its inverted page table occupies $2^{26}$ bytes,
  \SI{64}{\mebi\byte}, whether the machine runs ten processes or a thousand.
\end{example}

\section{Segmentation}
\label{sec:l08-segmentation}

With segmentation the address space is divided into segments of arbitrary
size that correspond to logically meaningful components of the process,
such as its code, data, heap and stack.\source{P3L2, segmentation; OSTEP
ch.~16; Silberschatz et al.\ ch.~9; Bovet and Cesati ch.~2.}  A virtual
address, on x86 called a \emph{logical address}, consists of a \emph{segment
selector} and an offset within the segment.  In its simplest form a segment
is a contiguous region of physical memory described by its \emph{base}, the
address at which it starts, and its \emph{limit}, its size.  The physical
address is $\text{base} + \text{offset}$, and it is valid only if
$\text{offset} < \text{limit}$; otherwise the access faults.

\begin{example}
  A process has a code segment with base \texttt{0x4000} and limit
  \texttt{0x1800}, and a stack segment with base \texttt{0x9000} and limit
  \texttt{0x0800}.  The logical address (code, \texttt{0x0A20}) translates
  to \texttt{0x4A20}, since $\texttt{0x0A20} < \texttt{0x1800}$.  The
  logical address (stack, \texttt{0x0900}) faults, since its offset exceeds
  the limit of the stack segment.
\end{example}

On x86 the selector is held in a segment register, such as CS for code, DS
for data and SS for the stack.  It indexes a \emph{segment descriptor
table}---the global descriptor table (GDT) or a local descriptor table
(LDT)---whose entry, the segment descriptor, holds the base, the limit and
the permissions of the segment.  The sum of base and offset is a
\emph{linear address}.  Segmentation is combined with paging: the linear
address is translated into a physical address by the page tables
(\cref{fig:l08-segmentation}).
\begin{itemize}
  \item On 32-bit x86 (IA-32) both mechanisms are active.  A descriptor
    table holds up to 8192 segment descriptors.  Linux, like most
    operating systems, defines segments that start at address~0 and span
    the whole address space, and it relies on paging for translation and
    protection.
  \item On x86-64, segmentation is retained only for backward
    compatibility.  In 64-bit mode the processor treats the segment bases as
    0 (except for FS and GS) and does not check limits, so paging alone
    translates addresses.
\end{itemize}

\begin{figure}[htbp]
  \centering
  % Segmentation combined with paging (x86): the selector indexes the segment
% descriptor table; the offset is checked against the limit and added to
% the base, giving a linear address that the page tables translate.
% Usage: % Segmentation combined with paging (x86): the selector indexes the segment
% descriptor table; the offset is checked against the limit and added to
% the base, giving a linear address that the page tables translate.
% Usage: % Segmentation combined with paging (x86): the selector indexes the segment
% descriptor table; the offset is checked against the limit and added to
% the base, giving a linear address that the page tables translate.
% Usage: \input{figures/l08-segmentation}   (width ~118 mm)
\begin{tikzpicture}[x=1mm, y=1mm, font=\scriptsize\sffamily,
  >={Stealth[length=1.5mm]},
  cell/.style={draw, minimum height=6mm, inner sep=0pt, anchor=south west},
  box/.style={draw, fill=white, align=center, inner sep=1.5pt, minimum height=9mm},
  lab/.style={font=\scriptsize\sffamily, align=center},
  hd/.style={font=\scriptsize\sffamily\bfseries, anchor=south west}]
  % logical address
  \node[hd] at (0,44.5) {\iflangja{論理アドレス}{logical address}};
  \node[cell, minimum width=16mm, fill=ln-source!16] at (0,38) {\iflangja{セレクタ}{selector}};
  \node[cell, minimum width=16mm, fill=ln-box] at (16,38) {\iflangja{オフセット}{offset}};
  % descriptor table
  \draw[fill=white] (0,28) rectangle (28,33);
  \draw[fill=ln-source!16] (0,22) rectangle (28,28) node[midway] {\iflangja{リミット}{limit}};
  \draw[fill=ln-source!16] (0,16) rectangle (28,22) node[midway] {\iflangja{ベース}{base}};
  \draw[fill=white] (0,11) rectangle (28,16);
  \draw[fill=white] (0,6) rectangle (28,11);
  \node[lab, anchor=north] at (14,5.5) {\iflangja{セグメント記述子表\\（GDT/LDT）}{segment descriptor\\table (GDT/LDT)}};
  \draw[->] (8,38) -- (8,33) node[pos=0.5, right, font=\scriptsize\sffamily, text=ln-muted] {\iflangja{索引}{index}};
  % limit check
  \node[box, minimum width=20mm] (chk) at (48,31) {\iflangja{オフセット\\$<$ リミット？}{offset\\$<$ limit?}};
  \draw[->] (32,41) -| (chk.north);
  \draw[->] (28,25) -- (34,25) |- (chk.west);
  \node[lab, anchor=west, text=ln-caution] (flt) at (64,31) {\iflangja{いいえ：フォールト}{no: fault}};
  \draw[->, ln-caution] (chk.east) -- (flt.west);
  % adder
  \node[draw, circle, inner sep=0.5pt, fill=white, font=\footnotesize] (add) at (48,13) {$+$};
  \draw[->] (chk.south) -- (add.north) node[pos=0.5, right, font=\scriptsize\sffamily] {\iflangja{はい}{yes}};
  \draw[->] (28,19) -| (40,13) -- (add.west);
  % linear -> paging -> physical
  \node[box, minimum width=18mm, fill=ln-accent!10] (lin) at (68,13) {\iflangja{リニア\\アドレス}{linear\\address}};
  \node[box, minimum width=18mm, fill=ln-accent!22] (pg) at (91,13) {\iflangja{ページ\\テーブル}{page\\tables}};
  \node[box, minimum width=16mm] (pa) at (111,13) {\iflangja{物理\\アドレス}{physical\\address}};
  \draw[->] (add.east) -- (lin.west);
  \draw[->] (lin.east) -- (pg.west);
  \draw[->] (pg.east) -- (pa.west);
\end{tikzpicture}
   (width ~118 mm)
\begin{tikzpicture}[x=1mm, y=1mm, font=\scriptsize\sffamily,
  >={Stealth[length=1.5mm]},
  cell/.style={draw, minimum height=6mm, inner sep=0pt, anchor=south west},
  box/.style={draw, fill=white, align=center, inner sep=1.5pt, minimum height=9mm},
  lab/.style={font=\scriptsize\sffamily, align=center},
  hd/.style={font=\scriptsize\sffamily\bfseries, anchor=south west}]
  % logical address
  \node[hd] at (0,44.5) {\iflangja{論理アドレス}{logical address}};
  \node[cell, minimum width=16mm, fill=ln-source!16] at (0,38) {\iflangja{セレクタ}{selector}};
  \node[cell, minimum width=16mm, fill=ln-box] at (16,38) {\iflangja{オフセット}{offset}};
  % descriptor table
  \draw[fill=white] (0,28) rectangle (28,33);
  \draw[fill=ln-source!16] (0,22) rectangle (28,28) node[midway] {\iflangja{リミット}{limit}};
  \draw[fill=ln-source!16] (0,16) rectangle (28,22) node[midway] {\iflangja{ベース}{base}};
  \draw[fill=white] (0,11) rectangle (28,16);
  \draw[fill=white] (0,6) rectangle (28,11);
  \node[lab, anchor=north] at (14,5.5) {\iflangja{セグメント記述子表\\（GDT/LDT）}{segment descriptor\\table (GDT/LDT)}};
  \draw[->] (8,38) -- (8,33) node[pos=0.5, right, font=\scriptsize\sffamily, text=ln-muted] {\iflangja{索引}{index}};
  % limit check
  \node[box, minimum width=20mm] (chk) at (48,31) {\iflangja{オフセット\\$<$ リミット？}{offset\\$<$ limit?}};
  \draw[->] (32,41) -| (chk.north);
  \draw[->] (28,25) -- (34,25) |- (chk.west);
  \node[lab, anchor=west, text=ln-caution] (flt) at (64,31) {\iflangja{いいえ：フォールト}{no: fault}};
  \draw[->, ln-caution] (chk.east) -- (flt.west);
  % adder
  \node[draw, circle, inner sep=0.5pt, fill=white, font=\footnotesize] (add) at (48,13) {$+$};
  \draw[->] (chk.south) -- (add.north) node[pos=0.5, right, font=\scriptsize\sffamily] {\iflangja{はい}{yes}};
  \draw[->] (28,19) -| (40,13) -- (add.west);
  % linear -> paging -> physical
  \node[box, minimum width=18mm, fill=ln-accent!10] (lin) at (68,13) {\iflangja{リニア\\アドレス}{linear\\address}};
  \node[box, minimum width=18mm, fill=ln-accent!22] (pg) at (91,13) {\iflangja{ページ\\テーブル}{page\\tables}};
  \node[box, minimum width=16mm] (pa) at (111,13) {\iflangja{物理\\アドレス}{physical\\address}};
  \draw[->] (add.east) -- (lin.west);
  \draw[->] (lin.east) -- (pg.west);
  \draw[->] (pg.east) -- (pa.west);
\end{tikzpicture}
   (width ~118 mm)
\begin{tikzpicture}[x=1mm, y=1mm, font=\scriptsize\sffamily,
  >={Stealth[length=1.5mm]},
  cell/.style={draw, minimum height=6mm, inner sep=0pt, anchor=south west},
  box/.style={draw, fill=white, align=center, inner sep=1.5pt, minimum height=9mm},
  lab/.style={font=\scriptsize\sffamily, align=center},
  hd/.style={font=\scriptsize\sffamily\bfseries, anchor=south west}]
  % logical address
  \node[hd] at (0,44.5) {\iflangja{論理アドレス}{logical address}};
  \node[cell, minimum width=16mm, fill=ln-source!16] at (0,38) {\iflangja{セレクタ}{selector}};
  \node[cell, minimum width=16mm, fill=ln-box] at (16,38) {\iflangja{オフセット}{offset}};
  % descriptor table
  \draw[fill=white] (0,28) rectangle (28,33);
  \draw[fill=ln-source!16] (0,22) rectangle (28,28) node[midway] {\iflangja{リミット}{limit}};
  \draw[fill=ln-source!16] (0,16) rectangle (28,22) node[midway] {\iflangja{ベース}{base}};
  \draw[fill=white] (0,11) rectangle (28,16);
  \draw[fill=white] (0,6) rectangle (28,11);
  \node[lab, anchor=north] at (14,5.5) {\iflangja{セグメント記述子表\\（GDT/LDT）}{segment descriptor\\table (GDT/LDT)}};
  \draw[->] (8,38) -- (8,33) node[pos=0.5, right, font=\scriptsize\sffamily, text=ln-muted] {\iflangja{索引}{index}};
  % limit check
  \node[box, minimum width=20mm] (chk) at (48,31) {\iflangja{オフセット\\$<$ リミット？}{offset\\$<$ limit?}};
  \draw[->] (32,41) -| (chk.north);
  \draw[->] (28,25) -- (34,25) |- (chk.west);
  \node[lab, anchor=west, text=ln-caution] (flt) at (64,31) {\iflangja{いいえ：フォールト}{no: fault}};
  \draw[->, ln-caution] (chk.east) -- (flt.west);
  % adder
  \node[draw, circle, inner sep=0.5pt, fill=white, font=\footnotesize] (add) at (48,13) {$+$};
  \draw[->] (chk.south) -- (add.north) node[pos=0.5, right, font=\scriptsize\sffamily] {\iflangja{はい}{yes}};
  \draw[->] (28,19) -| (40,13) -- (add.west);
  % linear -> paging -> physical
  \node[box, minimum width=18mm, fill=ln-accent!10] (lin) at (68,13) {\iflangja{リニア\\アドレス}{linear\\address}};
  \node[box, minimum width=18mm, fill=ln-accent!22] (pg) at (91,13) {\iflangja{ページ\\テーブル}{page\\tables}};
  \node[box, minimum width=16mm] (pa) at (111,13) {\iflangja{物理\\アドレス}{physical\\address}};
  \draw[->] (add.east) -- (lin.west);
  \draw[->] (lin.east) -- (pg.west);
  \draw[->] (pg.east) -- (pa.west);
\end{tikzpicture}

  \caption{Segmentation combined with paging on x86.  The selector indexes
    the segment descriptor table; an offset within the limit is added to the
    base, and the resulting linear address is translated by the page
    tables.}
  \label{fig:l08-segmentation}
\end{figure}

\section{Page size}
\label{sec:l08-page-size}

The size of a page is determined by the number of bits of its
offset.\source{P3L2, page size; OSTEP ch.~20, 23; Silberschatz et al.\
ch.~9.}  An offset of 10 bits gives pages of \SI{1}{\kibi\byte}, an offset
of 12 bits pages of \SI{4}{\kibi\byte}, which is the default on x86.  Many
architectures support several page sizes at the same time.

On x86, two larger page sizes are available besides the default.  A
\term{large page}[ラージページ] has a size of \SI{2}{\mebi\byte}, and its
offset has 21 bits.  A page of \SI{1}{\gibi\byte}, whose offset has 30
bits, is called a \term{huge page}[ヒュージページ].
\Cref{tab:l08-page-sizes} compares the three sizes.  (Linux calls pages of
both larger sizes huge pages.)

\begin{table}[htbp]
  \centering
  \caption{Page sizes on x86 and the page table entries needed to map
    \SI{1}{\gibi\byte} of memory.}
  \label{tab:l08-page-sizes}
  \small
  \begin{tabular}{@{}lrrr@{}}
    \toprule
    Page size & Offset bits & Entries for \SI{1}{\gibi\byte} & Reduction \\
    \midrule
    \SI{4}{\kibi\byte} (default) & 12 & 262\,144 & 1 \\
    \SI{2}{\mebi\byte} (large)   & 21 & 512      & $2^{9} = 512$ \\
    \SI{1}{\gibi\byte} (huge)    & 30 & 1        & $2^{18} = 262\,144$ \\
    \bottomrule
  \end{tabular}
\end{table}

Larger pages have two benefits.  They need fewer page table entries, so the
page tables are smaller; and each TLB entry covers more memory, so more
accesses are TLB hits.  Their drawback is wasted memory.

\begin{definition}[Internal fragmentation]
  Memory that is allocated but unused because an allocation is rounded up
  to a whole unit, here to a whole page, is lost to
  \term{internal fragmentation}[内部断片化].
\end{definition}

A large page is mapped as a whole even if the process uses only part of it,
so the larger the pages, the more memory internal fragmentation can
waste.

\begin{example}
  A process maps \SI{1}{\gibi\byte} of memory.  With \SI{4}{\kibi\byte}
  pages this takes 262\,144 entries, and a TLB of 64 entries reaches only
  \SI{256}{\kibi\byte} of it; with \SI{2}{\mebi\byte} pages it takes 512
  entries, and the same TLB reaches \SI{128}{\mebi\byte}.  A second region
  of \SI{5222}{\kibi\byte}, however, occupies 1306 pages of \SI{4}{\kibi\byte}, wasting
  \SI{2}{\kibi\byte}, but three pages of \SI{2}{\mebi\byte}, wasting
  $6144 - 5222 = \SI{922}{\kibi\byte}$.
\end{example}

\section{Memory allocation}
\label{sec:l08-allocation}

Address translation applies a mapping from virtual to physical
addresses.\source{P3L2, memory allocation, memory allocation challenges;
OSTEP ch.~17; Silberschatz et al.\ ch.~9--10.}  Page tables record the
mapping and the MMU checks every access against it, but some other
component must decide what the mapping is: which physical memory backs each
page of each process, and which memory is still free.

\begin{definition}[Memory allocator]
  A \term{memory allocator}[メモリアロケータ] determines which physical
  memory backs a region of virtual memory, and thereby the mapping from
  virtual to physical addresses.  Page tables and the MMU only record this
  mapping and check the validity of accesses and their permissions.
\end{definition}

Allocators exist at two levels.
\begin{description}
  \item[Kernel-level allocators] allocate memory regions, i.e., pages, for
    the state of the kernel, such as the process control blocks, and for the
    static state of processes: their code, stack and initialized data.  They
    also keep track of the free memory in the system.
  \item[User-level allocators] manage the dynamic state of a process, its
    heap, through calls such as \texttt{malloc} and \texttt{free}.  The
    kernel provides the process with a region of memory, and the user-level
    allocator divides it among the allocations without involving the kernel
    again.  Examples are dlmalloc, jemalloc, Hoard and tcmalloc.
\end{description}

The main challenge of an allocator is fragmentation of the free memory.

\begin{definition}[External fragmentation]
  Free memory suffers from \term{external fragmentation}[外部断片化] when
  it is split into non-contiguous holes, so that a request for a contiguous
  region cannot be satisfied although enough memory is free in total.
\end{definition}

Contiguous physical memory is needed, for example, for large pages, for
buffers that devices access directly, and for data structures of the
kernel.  An allocator counters external fragmentation in two ways: it
places allocations so that free memory stays contiguous, and when memory is
freed, it coalesces the freed region with adjacent free regions into a
larger one.

\begin{example}
  \label{ex:l08-fragmentation}
  A pool of 12 page frames receives requests for A (3 frames), B (2), C (3)
  and D (2); later A and C are freed, and a request for E needs 4
  contiguous frames (\cref{fig:l08-fragmentation}).  If A and C were placed
  apart, 8 frames are free, but in holes of 3, 3 and 2 frames: E fails.  If
  A and C were placed next to each other, freeing them leaves a contiguous
  hole of 6 frames, and E succeeds.
\end{example}

\begin{figure}[htbp]
  \centering
  % External fragmentation: 12 page frames, requests A(3) B(2) C(3) D(2),
% then A and C are freed and E(4) needs 4 contiguous frames.
% (a) A and C placed apart: E fails.  (b) A and C placed together: E fits.
% Usage: % External fragmentation: 12 page frames, requests A(3) B(2) C(3) D(2),
% then A and C are freed and E(4) needs 4 contiguous frames.
% (a) A and C placed apart: E fails.  (b) A and C placed together: E fits.
% Usage: % External fragmentation: 12 page frames, requests A(3) B(2) C(3) D(2),
% then A and C are freed and E(4) needs 4 contiguous frames.
% (a) A and C placed apart: E fails.  (b) A and C placed together: E fits.
% Usage: \input{figures/l08-fragmentation}   (width ~118 mm)
\begin{tikzpicture}[x=1mm, y=1mm, font=\scriptsize\sffamily,
  lab/.style={font=\scriptsize\sffamily, anchor=east},
  capt/.style={font=\footnotesize\sffamily, anchor=west},
  idx/.style={font=\scriptsize\sffamily, text=ln-muted}]
  % \frow{y}{list of 12 owners}: draw a row of 12 frames of 7.5 mm
  \def\frow#1#2{%
    \foreach \own [count=\c from 0] in {#2} {
      \if\own F\def\col{white}\def\txt{}\fi
      \if\own A\def\col{ln-accent!25}\def\txt{A}\fi
      \if\own B\def\col{ln-source!20}\def\txt{B}\fi
      \if\own C\def\col{ln-tip!22}\def\txt{C}\fi
      \if\own D\def\col{ln-caution!15}\def\txt{D}\fi
      \if\own E\def\col{ln-accent!55}\def\txt{E}\fi
      \draw[fill=\col] (28+7.5*\c,#1) rectangle ++(7.5,5)
        node[midway] {\txt};
    }}
  \foreach \c in {0,...,11} { \node[idx] at (31.75+7.5*\c,53) {\c}; }
  % ---------- (a)
  \node[capt] at (0,56) {\iflangja{(a) A と C を離して配置}{(a) A and C placed apart}};
  \node[lab] at (27,47.5) {\iflangja{A--D を割り当て後}{after A--D}};
  \frow{45}{A,A,A,B,B,C,C,C,D,D,F,F}
  \node[lab] at (27,40.5) {\iflangja{A，C を解放後}{after freeing A, C}};
  \frow{38}{F,F,F,B,B,F,F,F,D,D,F,F}
  \node[lab] at (27,33.5) {\iflangja{E（4 フレーム）}{request E (4)}};
  \node[font=\scriptsize\sffamily, text=ln-caution, anchor=west] at (28,33.5)
    {\iflangja{失敗：空き 8 フレームだが連続 4 フレームがない}{fails: 8 frames free, but no 4 contiguous}};
  % ---------- (b)
  \node[capt] at (0,24) {\iflangja{(b) A と C を隣接して配置}{(b) A and C placed together}};
  \node[lab] at (27,15.5) {\iflangja{A--D を割り当て後}{after A--D}};
  \frow{13}{A,A,A,C,C,C,B,B,D,D,F,F}
  \node[lab] at (27,8.5) {\iflangja{A，C を解放後}{after freeing A, C}};
  \frow{6}{F,F,F,F,F,F,B,B,D,D,F,F}
  \node[lab] at (27,1.5) {\iflangja{E（4 フレーム）}{request E (4)}};
  \frow{-1}{E,E,E,E,F,F,B,B,D,D,F,F}
\end{tikzpicture}
   (width ~118 mm)
\begin{tikzpicture}[x=1mm, y=1mm, font=\scriptsize\sffamily,
  lab/.style={font=\scriptsize\sffamily, anchor=east},
  capt/.style={font=\footnotesize\sffamily, anchor=west},
  idx/.style={font=\scriptsize\sffamily, text=ln-muted}]
  % \frow{y}{list of 12 owners}: draw a row of 12 frames of 7.5 mm
  \def\frow#1#2{%
    \foreach \own [count=\c from 0] in {#2} {
      \if\own F\def\col{white}\def\txt{}\fi
      \if\own A\def\col{ln-accent!25}\def\txt{A}\fi
      \if\own B\def\col{ln-source!20}\def\txt{B}\fi
      \if\own C\def\col{ln-tip!22}\def\txt{C}\fi
      \if\own D\def\col{ln-caution!15}\def\txt{D}\fi
      \if\own E\def\col{ln-accent!55}\def\txt{E}\fi
      \draw[fill=\col] (28+7.5*\c,#1) rectangle ++(7.5,5)
        node[midway] {\txt};
    }}
  \foreach \c in {0,...,11} { \node[idx] at (31.75+7.5*\c,53) {\c}; }
  % ---------- (a)
  \node[capt] at (0,56) {\iflangja{(a) A と C を離して配置}{(a) A and C placed apart}};
  \node[lab] at (27,47.5) {\iflangja{A--D を割り当て後}{after A--D}};
  \frow{45}{A,A,A,B,B,C,C,C,D,D,F,F}
  \node[lab] at (27,40.5) {\iflangja{A，C を解放後}{after freeing A, C}};
  \frow{38}{F,F,F,B,B,F,F,F,D,D,F,F}
  \node[lab] at (27,33.5) {\iflangja{E（4 フレーム）}{request E (4)}};
  \node[font=\scriptsize\sffamily, text=ln-caution, anchor=west] at (28,33.5)
    {\iflangja{失敗：空き 8 フレームだが連続 4 フレームがない}{fails: 8 frames free, but no 4 contiguous}};
  % ---------- (b)
  \node[capt] at (0,24) {\iflangja{(b) A と C を隣接して配置}{(b) A and C placed together}};
  \node[lab] at (27,15.5) {\iflangja{A--D を割り当て後}{after A--D}};
  \frow{13}{A,A,A,C,C,C,B,B,D,D,F,F}
  \node[lab] at (27,8.5) {\iflangja{A，C を解放後}{after freeing A, C}};
  \frow{6}{F,F,F,F,F,F,B,B,D,D,F,F}
  \node[lab] at (27,1.5) {\iflangja{E（4 フレーム）}{request E (4)}};
  \frow{-1}{E,E,E,E,F,F,B,B,D,D,F,F}
\end{tikzpicture}
   (width ~118 mm)
\begin{tikzpicture}[x=1mm, y=1mm, font=\scriptsize\sffamily,
  lab/.style={font=\scriptsize\sffamily, anchor=east},
  capt/.style={font=\footnotesize\sffamily, anchor=west},
  idx/.style={font=\scriptsize\sffamily, text=ln-muted}]
  % \frow{y}{list of 12 owners}: draw a row of 12 frames of 7.5 mm
  \def\frow#1#2{%
    \foreach \own [count=\c from 0] in {#2} {
      \if\own F\def\col{white}\def\txt{}\fi
      \if\own A\def\col{ln-accent!25}\def\txt{A}\fi
      \if\own B\def\col{ln-source!20}\def\txt{B}\fi
      \if\own C\def\col{ln-tip!22}\def\txt{C}\fi
      \if\own D\def\col{ln-caution!15}\def\txt{D}\fi
      \if\own E\def\col{ln-accent!55}\def\txt{E}\fi
      \draw[fill=\col] (28+7.5*\c,#1) rectangle ++(7.5,5)
        node[midway] {\txt};
    }}
  \foreach \c in {0,...,11} { \node[idx] at (31.75+7.5*\c,53) {\c}; }
  % ---------- (a)
  \node[capt] at (0,56) {\iflangja{(a) A と C を離して配置}{(a) A and C placed apart}};
  \node[lab] at (27,47.5) {\iflangja{A--D を割り当て後}{after A--D}};
  \frow{45}{A,A,A,B,B,C,C,C,D,D,F,F}
  \node[lab] at (27,40.5) {\iflangja{A，C を解放後}{after freeing A, C}};
  \frow{38}{F,F,F,B,B,F,F,F,D,D,F,F}
  \node[lab] at (27,33.5) {\iflangja{E（4 フレーム）}{request E (4)}};
  \node[font=\scriptsize\sffamily, text=ln-caution, anchor=west] at (28,33.5)
    {\iflangja{失敗：空き 8 フレームだが連続 4 フレームがない}{fails: 8 frames free, but no 4 contiguous}};
  % ---------- (b)
  \node[capt] at (0,24) {\iflangja{(b) A と C を隣接して配置}{(b) A and C placed together}};
  \node[lab] at (27,15.5) {\iflangja{A--D を割り当て後}{after A--D}};
  \frow{13}{A,A,A,C,C,C,B,B,D,D,F,F}
  \node[lab] at (27,8.5) {\iflangja{A，C を解放後}{after freeing A, C}};
  \frow{6}{F,F,F,F,F,F,B,B,D,D,F,F}
  \node[lab] at (27,1.5) {\iflangja{E（4 フレーム）}{request E (4)}};
  \frow{-1}{E,E,E,E,F,F,B,B,D,D,F,F}
\end{tikzpicture}

  \caption{External fragmentation in a pool of 12 frames.  The same requests
    and frees leave (a) three small holes or (b) one large hole, depending
    on where the allocator placed A and C.}
  \label{fig:l08-fragmentation}
\end{figure}

\section{Linux kernel allocators}
\label{sec:l08-linux-allocators}

The Linux kernel combines two allocators: a buddy allocator for page frames
and a slab allocator for kernel objects on top of it.\source{P3L2, Linux
kernel allocators; OSTEP ch.~17; Silberschatz et al.\ ch.~10; Bovet and
Cesati ch.~8; Bonwick \cite{bonwick1994}.}

\subsection{Buddy allocator}

\begin{definition}[Buddy allocator]
  A \term{buddy allocator}[バディアロケータ] manages free memory in blocks
  whose sizes are powers of two.  A request is rounded up to a power of two
  and served from the smallest free block that is large enough; a larger
  block is split in halves, called buddies, until a block of the requested
  size exists.  When a block is freed and its buddy is also free, the two
  are merged, and merging continues up the tree.
\end{definition}

Fragmentation can still occur, but merging works well and is fast: the
buddy of the block of $2^{k}$ frames that starts at frame $a$ starts at
frame $a \oplus 2^{k}$, i.e., the two addresses differ in a single bit, so
the allocator finds a buddy without searching.  The rounding to powers of
two, however, causes internal fragmentation.

\begin{example}
  \label{ex:l08-buddy}
  A buddy allocator manages 16 page frames.  \Cref{tab:l08-buddy} traces
  five operations, and \cref{fig:l08-buddy} shows the blocks after the
  third.  The request for 3 frames is rounded up to 4 and that for 7 frames
  to 8, wasting one frame each.  When B is freed, its buddy 6--7 is free,
  so the two merge into 4--7; the buddy of 4--7, block 0--3 with $0 = 4
  \oplus 4$, is also free, so the merge continues to 0--7, and it stops
  there because the buddy 8--15 is allocated.
\end{example}

\begin{table}[htbp]
  \centering
  \caption{Trace of a buddy allocator with 16 page frames.}
  \label{tab:l08-buddy}
  \small
  \begin{tabular}{@{}clll@{}}
    \toprule
    Step & Operation & Splits or merges & Free blocks after \\
    \midrule
    0 & ---                 & ---                                          & 0--15 \\
    1 & allocate A (3 $\to$ 4) & split 0--15, 0--7; A $=$ 0--3            & 4--7, 8--15 \\
    2 & allocate B (2)      & split 4--7; B $=$ 4--5                       & 6--7, 8--15 \\
    3 & allocate C (7 $\to$ 8) & C $=$ 8--15                              & 6--7 \\
    4 & free A              & buddy 4--7 is split: no merge                & 0--3, 6--7 \\
    5 & free B              & merge 4--5 with 6--7, then with 0--3         & 0--7 \\
    \bottomrule
  \end{tabular}
\end{table}

\begin{figure}[htbp]
  \centering
  % Buddy allocator over 16 page frames after A(3 -> 4), B(2) and C(7 -> 8)
% have been allocated.  Split blocks are dashed; the unused page of a
% rounded-up allocation is lightly shaded.
% Usage: % Buddy allocator over 16 page frames after A(3 -> 4), B(2) and C(7 -> 8)
% have been allocated.  Split blocks are dashed; the unused page of a
% rounded-up allocation is lightly shaded.
% Usage: % Buddy allocator over 16 page frames after A(3 -> 4), B(2) and C(7 -> 8)
% have been allocated.  Split blocks are dashed; the unused page of a
% rounded-up allocation is lightly shaded.
% Usage: \input{figures/l08-buddy}   (width ~118 mm)
\begin{tikzpicture}[x=1mm, y=1mm, font=\scriptsize\sffamily,
  split/.style={draw, densely dashed, fill=white},
  used/.style={draw, fill=ln-accent!25},
  waste/.style={draw, fill=ln-accent!8},
  freeb/.style={draw, fill=white},
  lab/.style={font=\scriptsize\sffamily, anchor=east, text=ln-muted},
  edge/.style={ln-rule}]
  % frame f starts at x = 22 + 6 f; levels at y = 33 (16), 22 (8), 11 (4), 0 (2)
  \node[lab] at (20,36) {\iflangja{16 フレーム}{16 frames}};
  \node[lab] at (20,25) {\iflangja{8 フレーム}{8 frames}};
  \node[lab] at (20,14) {\iflangja{4 フレーム}{4 frames}};
  \node[lab] at (20,3)  {\iflangja{2 フレーム}{2 frames}};
  % level 16
  \draw[split] (22,33) rectangle (118,39) node[midway] {0--15 \iflangja{（分割）}{(split)}};
  % level 8
  \draw[split] (22,22) rectangle (70,28) node[midway] {0--7 \iflangja{（分割）}{(split)}};
  \draw[used] (70,22) rectangle (112,28) node[midway] {C (7)};
  \draw[waste] (112,22) rectangle (118,28);
  % level 4
  \draw[used] (22,11) rectangle (40,17) node[midway] {A (3)};
  \draw[waste] (40,11) rectangle (46,17);
  \draw[split] (46,11) rectangle (70,17) node[midway] {4--7 \iflangja{（分割）}{(split)}};
  % level 2
  \draw[used] (46,0) rectangle (58,6) node[midway] {B};
  \draw[freeb] (58,0) rectangle (70,6) node[midway] {\iflangja{空き}{free}};
  % tree edges
  \draw[edge] (46,33) -- (46,28);
  \draw[edge] (94,33) -- (94,28);
  \draw[edge] (34,22) -- (34,17);
  \draw[edge] (58,22) -- (58,17);
  \draw[edge] (52,11) -- (52,6);
  \draw[edge] (64,11) -- (64,6);
  % frame axis
  \draw[ln-muted] (22,-3) -- (118,-3);
  \foreach \f in {0,4,8,12,16} {
    \draw[ln-muted] (22+6*\f,-3) -- ++(0,-1);
    \node[font=\scriptsize\sffamily, text=ln-muted, anchor=north] at (22+6*\f,-4) {\f};
  }
  \node[font=\scriptsize\sffamily, text=ln-muted, anchor=north west] at (22,-8)
    {\iflangja{フレーム番号}{frame number}};
\end{tikzpicture}
   (width ~118 mm)
\begin{tikzpicture}[x=1mm, y=1mm, font=\scriptsize\sffamily,
  split/.style={draw, densely dashed, fill=white},
  used/.style={draw, fill=ln-accent!25},
  waste/.style={draw, fill=ln-accent!8},
  freeb/.style={draw, fill=white},
  lab/.style={font=\scriptsize\sffamily, anchor=east, text=ln-muted},
  edge/.style={ln-rule}]
  % frame f starts at x = 22 + 6 f; levels at y = 33 (16), 22 (8), 11 (4), 0 (2)
  \node[lab] at (20,36) {\iflangja{16 フレーム}{16 frames}};
  \node[lab] at (20,25) {\iflangja{8 フレーム}{8 frames}};
  \node[lab] at (20,14) {\iflangja{4 フレーム}{4 frames}};
  \node[lab] at (20,3)  {\iflangja{2 フレーム}{2 frames}};
  % level 16
  \draw[split] (22,33) rectangle (118,39) node[midway] {0--15 \iflangja{（分割）}{(split)}};
  % level 8
  \draw[split] (22,22) rectangle (70,28) node[midway] {0--7 \iflangja{（分割）}{(split)}};
  \draw[used] (70,22) rectangle (112,28) node[midway] {C (7)};
  \draw[waste] (112,22) rectangle (118,28);
  % level 4
  \draw[used] (22,11) rectangle (40,17) node[midway] {A (3)};
  \draw[waste] (40,11) rectangle (46,17);
  \draw[split] (46,11) rectangle (70,17) node[midway] {4--7 \iflangja{（分割）}{(split)}};
  % level 2
  \draw[used] (46,0) rectangle (58,6) node[midway] {B};
  \draw[freeb] (58,0) rectangle (70,6) node[midway] {\iflangja{空き}{free}};
  % tree edges
  \draw[edge] (46,33) -- (46,28);
  \draw[edge] (94,33) -- (94,28);
  \draw[edge] (34,22) -- (34,17);
  \draw[edge] (58,22) -- (58,17);
  \draw[edge] (52,11) -- (52,6);
  \draw[edge] (64,11) -- (64,6);
  % frame axis
  \draw[ln-muted] (22,-3) -- (118,-3);
  \foreach \f in {0,4,8,12,16} {
    \draw[ln-muted] (22+6*\f,-3) -- ++(0,-1);
    \node[font=\scriptsize\sffamily, text=ln-muted, anchor=north] at (22+6*\f,-4) {\f};
  }
  \node[font=\scriptsize\sffamily, text=ln-muted, anchor=north west] at (22,-8)
    {\iflangja{フレーム番号}{frame number}};
\end{tikzpicture}
   (width ~118 mm)
\begin{tikzpicture}[x=1mm, y=1mm, font=\scriptsize\sffamily,
  split/.style={draw, densely dashed, fill=white},
  used/.style={draw, fill=ln-accent!25},
  waste/.style={draw, fill=ln-accent!8},
  freeb/.style={draw, fill=white},
  lab/.style={font=\scriptsize\sffamily, anchor=east, text=ln-muted},
  edge/.style={ln-rule}]
  % frame f starts at x = 22 + 6 f; levels at y = 33 (16), 22 (8), 11 (4), 0 (2)
  \node[lab] at (20,36) {\iflangja{16 フレーム}{16 frames}};
  \node[lab] at (20,25) {\iflangja{8 フレーム}{8 frames}};
  \node[lab] at (20,14) {\iflangja{4 フレーム}{4 frames}};
  \node[lab] at (20,3)  {\iflangja{2 フレーム}{2 frames}};
  % level 16
  \draw[split] (22,33) rectangle (118,39) node[midway] {0--15 \iflangja{（分割）}{(split)}};
  % level 8
  \draw[split] (22,22) rectangle (70,28) node[midway] {0--7 \iflangja{（分割）}{(split)}};
  \draw[used] (70,22) rectangle (112,28) node[midway] {C (7)};
  \draw[waste] (112,22) rectangle (118,28);
  % level 4
  \draw[used] (22,11) rectangle (40,17) node[midway] {A (3)};
  \draw[waste] (40,11) rectangle (46,17);
  \draw[split] (46,11) rectangle (70,17) node[midway] {4--7 \iflangja{（分割）}{(split)}};
  % level 2
  \draw[used] (46,0) rectangle (58,6) node[midway] {B};
  \draw[freeb] (58,0) rectangle (70,6) node[midway] {\iflangja{空き}{free}};
  % tree edges
  \draw[edge] (46,33) -- (46,28);
  \draw[edge] (94,33) -- (94,28);
  \draw[edge] (34,22) -- (34,17);
  \draw[edge] (58,22) -- (58,17);
  \draw[edge] (52,11) -- (52,6);
  \draw[edge] (64,11) -- (64,6);
  % frame axis
  \draw[ln-muted] (22,-3) -- (118,-3);
  \foreach \f in {0,4,8,12,16} {
    \draw[ln-muted] (22+6*\f,-3) -- ++(0,-1);
    \node[font=\scriptsize\sffamily, text=ln-muted, anchor=north] at (22+6*\f,-4) {\f};
  }
  \node[font=\scriptsize\sffamily, text=ln-muted, anchor=north west] at (22,-8)
    {\iflangja{フレーム番号}{frame number}};
\end{tikzpicture}

  \caption{The blocks of the buddy allocator of \cref{ex:l08-buddy} after
    step~3, with the number of frames requested in parentheses.  Dashed
    blocks are split; the lightly shaded frames are lost to internal
    fragmentation.}
  \label{fig:l08-buddy}
\end{figure}

\subsection{Slab allocator}

Objects of the kernel have sizes that are not powers of two, and the buddy
allocator would waste much of the memory it gives them.  The Linux kernel
therefore allocates its objects with a slab allocator \cite{bonwick1994}.

\begin{definition}[Slab allocator]
  A \term{slab allocator}[スラブアロケータ] keeps a cache for every common
  object type or size.  Each cache consists of \emph{slabs}, regions of
  contiguously allocated physical memory, divided into slots of exactly the
  size of the object.  An allocation takes a free slot from the cache of
  the object's type, and a free returns the slot to the cache.
\end{definition}

Since slots have exactly the size of the objects, the slab allocator avoids
the rounding waste of the buddy allocator---only the remainder at the end of
a slab is unused---and external fragmentation is not an issue: a slot that is
freed fits the next object of the same type (\cref{fig:l08-slab}).  The
underlying slabs are obtained from the buddy allocator.

\begin{figure}[htbp]
  \centering
  % Slab allocator: one cache per object type; each cache consists of slabs
% (contiguous page frames) divided into slots of exactly the object size.
% Usage: % Slab allocator: one cache per object type; each cache consists of slabs
% (contiguous page frames) divided into slots of exactly the object size.
% Usage: % Slab allocator: one cache per object type; each cache consists of slabs
% (contiguous page frames) divided into slots of exactly the object size.
% Usage: \input{figures/l08-slab}   (width ~118 mm)
\begin{tikzpicture}[x=1mm, y=1mm, font=\scriptsize\sffamily,
  usedslot/.style={draw, fill=ln-accent!25},
  freeslot/.style={draw, fill=white},
  cachebox/.style={draw, fill=ln-source!14, align=center, minimum width=24mm,
                   minimum height=10mm, inner sep=1pt},
  lab/.style={font=\scriptsize\sffamily, text=ln-muted},
  >={Stealth[length=1.4mm]}]
  % cache 1: small objects, 8 slots per slab
  \node[cachebox] (c1) at (12,27) {\iflangja{キャッシュ 1\\（小さなオブジェクト）}{cache 1\\(small objects)}};
  \foreach \k/\s in {0/usedslot,1/usedslot,2/freeslot,3/usedslot,4/usedslot,5/usedslot,6/freeslot,7/usedslot} {
    \draw[\s] (32+5*\k,24) rectangle ++(5,6);
  }
  \foreach \k/\s in {0/usedslot,1/freeslot,2/freeslot,3/freeslot,4/freeslot,5/freeslot,6/freeslot,7/freeslot} {
    \draw[\s] (78+5*\k,24) rectangle ++(5,6);
  }
  \draw[thick] (32,24) rectangle (72,30);
  \draw[thick] (78,24) rectangle (118,30);
  \draw[->] (c1.east) -- (32,27);
  \draw[->] (72,27) -- (78,27);
  % cache 2: larger objects, 3 slots per slab
  \node[cachebox] (c2) at (12,9) {\iflangja{キャッシュ 2\\（大きなオブジェクト）}{cache 2\\(large objects)}};
  \foreach \k/\s in {0/usedslot,1/usedslot,2/usedslot} {
    \draw[\s] (32+13.33*\k,6) rectangle ++(13.33,6);
  }
  \foreach \k/\s in {0/freeslot,1/usedslot,2/freeslot} {
    \draw[\s] (78+13.33*\k,6) rectangle ++(13.33,6);
  }
  \draw[thick] (32,6) rectangle (72,12);
  \draw[thick] (78,6) rectangle (118,12);
  \draw[->] (c2.east) -- (32,9);
  \draw[->] (72,9) -- (78,9);
  % labels
  \node[lab, anchor=south] at (52,30.5) {\iflangja{スラブ（連続したページフレーム）}{slab (contiguous page frames)}};
  \node[lab, anchor=south] at (98,30.5) {\iflangja{スラブ}{slab}};
  \node[lab, anchor=north] at (52,5.3) {\iflangja{スロット＝オブジェクトの大きさ}{slot size = object size}};
  \draw[usedslot] (76,-0.5) rectangle ++(3,3);
  \node[lab, anchor=west] at (79.5,1) {\iflangja{使用中}{allocated}};
  \draw[freeslot] (100,-0.5) rectangle ++(3,3);
  \node[lab, anchor=west] at (103.5,1) {\iflangja{空き}{free}};
\end{tikzpicture}
   (width ~118 mm)
\begin{tikzpicture}[x=1mm, y=1mm, font=\scriptsize\sffamily,
  usedslot/.style={draw, fill=ln-accent!25},
  freeslot/.style={draw, fill=white},
  cachebox/.style={draw, fill=ln-source!14, align=center, minimum width=24mm,
                   minimum height=10mm, inner sep=1pt},
  lab/.style={font=\scriptsize\sffamily, text=ln-muted},
  >={Stealth[length=1.4mm]}]
  % cache 1: small objects, 8 slots per slab
  \node[cachebox] (c1) at (12,27) {\iflangja{キャッシュ 1\\（小さなオブジェクト）}{cache 1\\(small objects)}};
  \foreach \k/\s in {0/usedslot,1/usedslot,2/freeslot,3/usedslot,4/usedslot,5/usedslot,6/freeslot,7/usedslot} {
    \draw[\s] (32+5*\k,24) rectangle ++(5,6);
  }
  \foreach \k/\s in {0/usedslot,1/freeslot,2/freeslot,3/freeslot,4/freeslot,5/freeslot,6/freeslot,7/freeslot} {
    \draw[\s] (78+5*\k,24) rectangle ++(5,6);
  }
  \draw[thick] (32,24) rectangle (72,30);
  \draw[thick] (78,24) rectangle (118,30);
  \draw[->] (c1.east) -- (32,27);
  \draw[->] (72,27) -- (78,27);
  % cache 2: larger objects, 3 slots per slab
  \node[cachebox] (c2) at (12,9) {\iflangja{キャッシュ 2\\（大きなオブジェクト）}{cache 2\\(large objects)}};
  \foreach \k/\s in {0/usedslot,1/usedslot,2/usedslot} {
    \draw[\s] (32+13.33*\k,6) rectangle ++(13.33,6);
  }
  \foreach \k/\s in {0/freeslot,1/usedslot,2/freeslot} {
    \draw[\s] (78+13.33*\k,6) rectangle ++(13.33,6);
  }
  \draw[thick] (32,6) rectangle (72,12);
  \draw[thick] (78,6) rectangle (118,12);
  \draw[->] (c2.east) -- (32,9);
  \draw[->] (72,9) -- (78,9);
  % labels
  \node[lab, anchor=south] at (52,30.5) {\iflangja{スラブ（連続したページフレーム）}{slab (contiguous page frames)}};
  \node[lab, anchor=south] at (98,30.5) {\iflangja{スラブ}{slab}};
  \node[lab, anchor=north] at (52,5.3) {\iflangja{スロット＝オブジェクトの大きさ}{slot size = object size}};
  \draw[usedslot] (76,-0.5) rectangle ++(3,3);
  \node[lab, anchor=west] at (79.5,1) {\iflangja{使用中}{allocated}};
  \draw[freeslot] (100,-0.5) rectangle ++(3,3);
  \node[lab, anchor=west] at (103.5,1) {\iflangja{空き}{free}};
\end{tikzpicture}
   (width ~118 mm)
\begin{tikzpicture}[x=1mm, y=1mm, font=\scriptsize\sffamily,
  usedslot/.style={draw, fill=ln-accent!25},
  freeslot/.style={draw, fill=white},
  cachebox/.style={draw, fill=ln-source!14, align=center, minimum width=24mm,
                   minimum height=10mm, inner sep=1pt},
  lab/.style={font=\scriptsize\sffamily, text=ln-muted},
  >={Stealth[length=1.4mm]}]
  % cache 1: small objects, 8 slots per slab
  \node[cachebox] (c1) at (12,27) {\iflangja{キャッシュ 1\\（小さなオブジェクト）}{cache 1\\(small objects)}};
  \foreach \k/\s in {0/usedslot,1/usedslot,2/freeslot,3/usedslot,4/usedslot,5/usedslot,6/freeslot,7/usedslot} {
    \draw[\s] (32+5*\k,24) rectangle ++(5,6);
  }
  \foreach \k/\s in {0/usedslot,1/freeslot,2/freeslot,3/freeslot,4/freeslot,5/freeslot,6/freeslot,7/freeslot} {
    \draw[\s] (78+5*\k,24) rectangle ++(5,6);
  }
  \draw[thick] (32,24) rectangle (72,30);
  \draw[thick] (78,24) rectangle (118,30);
  \draw[->] (c1.east) -- (32,27);
  \draw[->] (72,27) -- (78,27);
  % cache 2: larger objects, 3 slots per slab
  \node[cachebox] (c2) at (12,9) {\iflangja{キャッシュ 2\\（大きなオブジェクト）}{cache 2\\(large objects)}};
  \foreach \k/\s in {0/usedslot,1/usedslot,2/usedslot} {
    \draw[\s] (32+13.33*\k,6) rectangle ++(13.33,6);
  }
  \foreach \k/\s in {0/freeslot,1/usedslot,2/freeslot} {
    \draw[\s] (78+13.33*\k,6) rectangle ++(13.33,6);
  }
  \draw[thick] (32,6) rectangle (72,12);
  \draw[thick] (78,6) rectangle (118,12);
  \draw[->] (c2.east) -- (32,9);
  \draw[->] (72,9) -- (78,9);
  % labels
  \node[lab, anchor=south] at (52,30.5) {\iflangja{スラブ（連続したページフレーム）}{slab (contiguous page frames)}};
  \node[lab, anchor=south] at (98,30.5) {\iflangja{スラブ}{slab}};
  \node[lab, anchor=north] at (52,5.3) {\iflangja{スロット＝オブジェクトの大きさ}{slot size = object size}};
  \draw[usedslot] (76,-0.5) rectangle ++(3,3);
  \node[lab, anchor=west] at (79.5,1) {\iflangja{使用中}{allocated}};
  \draw[freeslot] (100,-0.5) rectangle ++(3,3);
  \node[lab, anchor=west] at (103.5,1) {\iflangja{空き}{free}};
\end{tikzpicture}

  \caption{A slab allocator with a cache for small and a cache for large
    objects.  Every slab is divided into slots of the object size.}
  \label{fig:l08-slab}
\end{figure}

\begin{example}
  A kernel object occupies 1200 bytes.  Allocated separately from a buddy
  allocator, each object would receive at least one page of
  \SI{4}{\kibi\byte} and waste 2896 bytes.  A slab of one page holds
  $\lfloor 4096/1200 \rfloor = 3$ such objects and leaves $4096 - 3600 =
  496$ bytes unused, about 165 bytes per object.
\end{example}

\needspace{9\baselineskip}
\begin{keypoint}
  Linux allocates page frames with a buddy allocator, whose power-of-two
  blocks are found and merged without searching but which rounds requests
  up, and kernel objects with a slab allocator, whose slots have the size
  of the object and are reused by the next object of the same type.
\end{keypoint}

\section{Demand paging}
\label{sec:l08-demand}

Since virtual memory is far larger than physical memory, a virtual page is
not always present in a page frame.\source{P3L2, demand paging; OSTEP
ch.~21; Silberschatz et al.\ ch.~10.}  The contents of page frames are
saved to and restored from secondary storage.

\begin{definition}[Swapping and demand paging]
  Moving the contents of pages between physical memory and secondary
  storage is called \term{swapping}[スワッピング].  The pages are stored in
  a \term{swap partition}[スワップ領域], a region of a disk, of flash
  memory, or even of the memory of a remote machine.  With
  \term{demand paging}[デマンドページング], a page is swapped into memory
  only when the process accesses it.
\end{definition}

When a process accesses a page that is not in memory, the following steps
take place (\cref{fig:l08-page-fault}):
\begin{enumerate}
  \item The process references a memory address, and the MMU finds the
    valid bit of its PTE cleared.
  \item The MMU raises a page fault, which traps into the operating system.
  \item The operating system determines that the access is legal and where
    the page is stored in the swap partition.
  \item It reads the page from the swap partition into a free page frame,
    using an I/O operation.
  \item It updates the PTE with the new PFN and sets the valid bit.
  \item It resets the program counter to the faulting instruction and lets
    the process continue, so that the reference is executed again and now
    succeeds.
\end{enumerate}
The page frame that the page occupies after it is swapped in is generally
not the one it occupied before.  If no page frame is free, the operating
system must first free one (\cref{sec:l08-replacement}).

\begin{figure}[htbp]
  \centering
  % Demand paging: the steps taken when a process references a page that is
% in the swap partition.
% Usage: % Demand paging: the steps taken when a process references a page that is
% in the swap partition.
% Usage: % Demand paging: the steps taken when a process references a page that is
% in the swap partition.
% Usage: \input{figures/l08-page-fault}   (width ~118 mm)
\begin{tikzpicture}[x=1mm, y=1mm, font=\scriptsize\sffamily,
  >={Stealth[length=1.5mm]},
  box/.style={draw, fill=white, align=center, inner sep=1pt},
  stp/.style={draw, circle, fill=ln-accent, text=white, inner sep=0.4pt,
               font=\scriptsize\sffamily\bfseries, minimum size=4mm},
  lab/.style={font=\scriptsize\sffamily, align=left},
  hd/.style={font=\scriptsize\sffamily\bfseries, anchor=south}]
  % process
  \node[box, minimum width=22mm, minimum height=13mm] (proc) at (11,38.5)
    {\iflangja{プロセス}{process}\\[1pt]\texttt{x = a[i];}};
  % page table
  \node[hd, anchor=north] at (46,25.5) {\iflangja{ページテーブル}{page table}};
  \foreach \r in {0,...,4} { \draw[fill=white] (38,26+5*\r) rectangle ++(16,5); }
  \draw[fill=ln-caution!15] (38,36) rectangle ++(16,5) node[midway] {\iflangja{無効}{invalid}};
  % OS
  \node[box, minimum width=44mm, minimum height=8mm, fill=ln-source!12] (os) at (46,62)
    {\iflangja{OS（ページフォールトハンドラ）}{OS (page fault handler)}};
  % physical memory
  \node[hd] at (104,50.5) {\iflangja{物理メモリ}{physical memory}};
  \foreach \r in {0,...,4} { \draw[fill=ln-box] (92,26+5*\r) rectangle ++(24,5); }
  \draw[fill=ln-accent!25] (92,26) rectangle ++(24,5) node[midway] {\iflangja{空きフレーム}{free frame}};
  % swap partition
  \node[box, minimum width=28mm, minimum height=8mm, fill=ln-box, rounded corners=3pt] (swap) at (104,6)
    {\iflangja{スワップ領域}{swap partition}};
  % 1 reference
  \draw[->] (proc.east) -- (38,38.5);
  \node[stp] at (30,41) {1};
  % 2 trap
  \draw[->] (42,51) -- (42,58);
  \node[stp] at (39,54.5) {2};
  \node[lab, anchor=east] at (37,54.5) {\iflangja{トラップ}{trap}};
  % 3 locate page in swap
  \draw[->] (68.5,62) -- (80,62) |- (swap.west);
  \node[stp] at (77,58) {3};
  \node[lab, anchor=east] at (79,46) {\iflangja{スワップ上の\\位置を特定}{locate page\\in swap}};
  % 4 read page into free frame
  \draw[->, very thick, ln-accent] (104,10) -- (104,26);
  \node[stp] at (107.5,17) {4};
  \node[lab, anchor=east] at (101,17) {\iflangja{読み込み}{read page in}};
  % 5 update page table entry
  \draw[->] (50,58) -- (50,51);
  \node[stp] at (53,54.5) {5};
  \node[lab, anchor=west] at (55,54.5) {\iflangja{PTE を更新}{update PTE}};
  \draw[->, densely dashed] (54,38.5) -- (92,28.5);
  % 6 restart
  \draw[->] (os.west) -| (11,44.5);
  \node[stp] at (14,58) {6};
  \node[lab, anchor=west] at (16,54) {\iflangja{命令を\\再実行}{restart}};
\end{tikzpicture}
   (width ~118 mm)
\begin{tikzpicture}[x=1mm, y=1mm, font=\scriptsize\sffamily,
  >={Stealth[length=1.5mm]},
  box/.style={draw, fill=white, align=center, inner sep=1pt},
  stp/.style={draw, circle, fill=ln-accent, text=white, inner sep=0.4pt,
               font=\scriptsize\sffamily\bfseries, minimum size=4mm},
  lab/.style={font=\scriptsize\sffamily, align=left},
  hd/.style={font=\scriptsize\sffamily\bfseries, anchor=south}]
  % process
  \node[box, minimum width=22mm, minimum height=13mm] (proc) at (11,38.5)
    {\iflangja{プロセス}{process}\\[1pt]\texttt{x = a[i];}};
  % page table
  \node[hd, anchor=north] at (46,25.5) {\iflangja{ページテーブル}{page table}};
  \foreach \r in {0,...,4} { \draw[fill=white] (38,26+5*\r) rectangle ++(16,5); }
  \draw[fill=ln-caution!15] (38,36) rectangle ++(16,5) node[midway] {\iflangja{無効}{invalid}};
  % OS
  \node[box, minimum width=44mm, minimum height=8mm, fill=ln-source!12] (os) at (46,62)
    {\iflangja{OS（ページフォールトハンドラ）}{OS (page fault handler)}};
  % physical memory
  \node[hd] at (104,50.5) {\iflangja{物理メモリ}{physical memory}};
  \foreach \r in {0,...,4} { \draw[fill=ln-box] (92,26+5*\r) rectangle ++(24,5); }
  \draw[fill=ln-accent!25] (92,26) rectangle ++(24,5) node[midway] {\iflangja{空きフレーム}{free frame}};
  % swap partition
  \node[box, minimum width=28mm, minimum height=8mm, fill=ln-box, rounded corners=3pt] (swap) at (104,6)
    {\iflangja{スワップ領域}{swap partition}};
  % 1 reference
  \draw[->] (proc.east) -- (38,38.5);
  \node[stp] at (30,41) {1};
  % 2 trap
  \draw[->] (42,51) -- (42,58);
  \node[stp] at (39,54.5) {2};
  \node[lab, anchor=east] at (37,54.5) {\iflangja{トラップ}{trap}};
  % 3 locate page in swap
  \draw[->] (68.5,62) -- (80,62) |- (swap.west);
  \node[stp] at (77,58) {3};
  \node[lab, anchor=east] at (79,46) {\iflangja{スワップ上の\\位置を特定}{locate page\\in swap}};
  % 4 read page into free frame
  \draw[->, very thick, ln-accent] (104,10) -- (104,26);
  \node[stp] at (107.5,17) {4};
  \node[lab, anchor=east] at (101,17) {\iflangja{読み込み}{read page in}};
  % 5 update page table entry
  \draw[->] (50,58) -- (50,51);
  \node[stp] at (53,54.5) {5};
  \node[lab, anchor=west] at (55,54.5) {\iflangja{PTE を更新}{update PTE}};
  \draw[->, densely dashed] (54,38.5) -- (92,28.5);
  % 6 restart
  \draw[->] (os.west) -| (11,44.5);
  \node[stp] at (14,58) {6};
  \node[lab, anchor=west] at (16,54) {\iflangja{命令を\\再実行}{restart}};
\end{tikzpicture}
   (width ~118 mm)
\begin{tikzpicture}[x=1mm, y=1mm, font=\scriptsize\sffamily,
  >={Stealth[length=1.5mm]},
  box/.style={draw, fill=white, align=center, inner sep=1pt},
  stp/.style={draw, circle, fill=ln-accent, text=white, inner sep=0.4pt,
               font=\scriptsize\sffamily\bfseries, minimum size=4mm},
  lab/.style={font=\scriptsize\sffamily, align=left},
  hd/.style={font=\scriptsize\sffamily\bfseries, anchor=south}]
  % process
  \node[box, minimum width=22mm, minimum height=13mm] (proc) at (11,38.5)
    {\iflangja{プロセス}{process}\\[1pt]\texttt{x = a[i];}};
  % page table
  \node[hd, anchor=north] at (46,25.5) {\iflangja{ページテーブル}{page table}};
  \foreach \r in {0,...,4} { \draw[fill=white] (38,26+5*\r) rectangle ++(16,5); }
  \draw[fill=ln-caution!15] (38,36) rectangle ++(16,5) node[midway] {\iflangja{無効}{invalid}};
  % OS
  \node[box, minimum width=44mm, minimum height=8mm, fill=ln-source!12] (os) at (46,62)
    {\iflangja{OS（ページフォールトハンドラ）}{OS (page fault handler)}};
  % physical memory
  \node[hd] at (104,50.5) {\iflangja{物理メモリ}{physical memory}};
  \foreach \r in {0,...,4} { \draw[fill=ln-box] (92,26+5*\r) rectangle ++(24,5); }
  \draw[fill=ln-accent!25] (92,26) rectangle ++(24,5) node[midway] {\iflangja{空きフレーム}{free frame}};
  % swap partition
  \node[box, minimum width=28mm, minimum height=8mm, fill=ln-box, rounded corners=3pt] (swap) at (104,6)
    {\iflangja{スワップ領域}{swap partition}};
  % 1 reference
  \draw[->] (proc.east) -- (38,38.5);
  \node[stp] at (30,41) {1};
  % 2 trap
  \draw[->] (42,51) -- (42,58);
  \node[stp] at (39,54.5) {2};
  \node[lab, anchor=east] at (37,54.5) {\iflangja{トラップ}{trap}};
  % 3 locate page in swap
  \draw[->] (68.5,62) -- (80,62) |- (swap.west);
  \node[stp] at (77,58) {3};
  \node[lab, anchor=east] at (79,46) {\iflangja{スワップ上の\\位置を特定}{locate page\\in swap}};
  % 4 read page into free frame
  \draw[->, very thick, ln-accent] (104,10) -- (104,26);
  \node[stp] at (107.5,17) {4};
  \node[lab, anchor=east] at (101,17) {\iflangja{読み込み}{read page in}};
  % 5 update page table entry
  \draw[->] (50,58) -- (50,51);
  \node[stp] at (53,54.5) {5};
  \node[lab, anchor=west] at (55,54.5) {\iflangja{PTE を更新}{update PTE}};
  \draw[->, densely dashed] (54,38.5) -- (92,28.5);
  % 6 restart
  \draw[->] (os.west) -| (11,44.5);
  \node[stp] at (14,58) {6};
  \node[lab, anchor=west] at (16,54) {\iflangja{命令を\\再実行}{restart}};
\end{tikzpicture}

  \caption{Demand paging: the steps taken when a process references a page
    that is in the swap partition.}
  \label{fig:l08-page-fault}
\end{figure}

Some pages must not be swapped.  A \term{pinned page}[ピン留めページ] is a
page for which swapping is disabled, so that it stays in physical memory.
Pages are pinned, for example, when a device accesses them directly by
their physical address during direct memory access (Lesson~11).

\begin{example}
  A memory access costs \SI{100}{\nano\second}, and handling a page fault
  that reads the page from disk costs \SI{8}{\milli\second}, i.e.,
  $8 \cdot 10^{6}$ nanoseconds, 80\,000 times as much.  With a fraction $p$
  of accesses faulting, the
  average access costs $(1-p) \cdot 100 + p \cdot 8 \cdot 10^{6}$
  nanoseconds.  For $p = 10^{-4}$ this is about \SI{900}{\nano\second}, nine
  times slower; to stay within \SI{10}{\percent} of \SI{100}{\nano\second},
  $p$ must not exceed $1.25 \cdot 10^{-6}$, one fault in 800\,000 accesses.
\end{example}

\section{Page replacement}
\label{sec:l08-replacement}

When physical memory runs short, the operating system must free page frames
by swapping pages out.\source{P3L2, freeing up physical memory; OSTEP
ch.~22; Silberschatz et al.\ ch.~10; Bovet and Cesati ch.~17.}  A page that
is swapped out but needed again soon costs a page fault and an I/O
operation, which is expensive, as the example of \cref{sec:l08-demand}
showed.  The operating system must therefore choose with care the pages
that it swaps out.  This choice is called
\term{page replacement}[ページ置換], and it raises two questions: when pages
should be swapped out, and which ones.

Pages are swapped out by a \term{page-out daemon}[ページアウトデーモン], a
kernel thread that runs when memory usage exceeds a threshold, the
\emph{high watermark}.  It preferably runs while CPU usage is below a
threshold, the \emph{low watermark}, so that it disturbs the applications
as little as possible.

The ideal victims are pages that will not be used in the future.  The
future is unknown, so the operating system predicts it from the history of
accesses.
\begin{itemize}
  \item The \term{least recently used}[LRU]\index{LRU|see{least recently used}}
    (LRU) policy swaps out the page whose last access lies furthest in the
    past, assuming that a page that has not been used recently will not be
    used soon.  The access bit of the PTE tells the operating system
    whether a page has been referenced since the bit was last cleared.
  \item Pages that need not be written out are cheaper to replace.  The
    dirty bit tells whether a page has been modified; the contents of a
    clean page that is already stored on disk can simply be discarded.
  \item Pages that must stay in memory, such as pinned pages and certain
    pages of the kernel, are never chosen.
\end{itemize}
Linux provides parameters that tune when reclaim starts and stops,
categorizes pages into types, for instance reclaimable and swappable pages,
and uses a variation of LRU that gives pages a second
chance.\caution{In Linux both watermarks are thresholds on free memory:
reclaim starts below the low and stops above the high one.}

\begin{definition}[Second chance algorithm]
  The \term{second chance algorithm}[セカンドチャンスアルゴリズム]
  approximates LRU with the access bit.  It scans the page frames in
  circular order.  A page whose access bit is set is given a second chance:
  its bit is cleared, and the scan moves on.  A page whose access bit is
  clear, i.e., one that has not been referenced since the previous scan
  cleared its bit, is swapped out.
\end{definition}

\begin{example}
  \label{ex:l08-second-chance}
  Five page frames hold pages A--E with access bits 1, 0, 1, 1, 0, and the
  scan starts at A.  \Cref{tab:l08-second-chance} traces two replacements.
  The second scan resumes after the frame that received F, so it examines
  only C, D and E: C and D are given a second chance, and E, whose bit is
  clear, is swapped out.  A is not examined at all; because the process
  referenced it again after the first scan had cleared its bit, it is given
  a second chance when the scan next reaches it.  C, whose bit was cleared
  in step~3 and which has not been referenced since, is then the victim.
\end{example}

\begin{table}[htbp]
  \centering
  \caption{Trace of the second chance algorithm.  Access bits are listed in
    scan order; a newly loaded page has its bit set.}
  \label{tab:l08-second-chance}
  \small
  \begin{tabular}{@{}clll@{}}
    \toprule
    Step & Event & Pages examined & Access bits after \\
    \midrule
    0 & ---                     & ---                          & A1 B0 C1 D1 E0 \\
    1 & page fault, load F      & A: clear; B: swap out        & A0 F1 C1 D1 E0 \\
    2 & process references A, D & ---                          & A1 F1 C1 D1 E0 \\
    3 & page fault, load G      & C: clear; D: clear; E: swap out & A1 F1 C0 D0 G1 \\
    \bottomrule
  \end{tabular}
\end{table}

\needspace{10\baselineskip}
\begin{keypoint}
  Demand paging brings a page in on the page fault that first needs it, and
  a page-out daemon frees frames before memory runs out.  Which pages leave
  is decided with an approximation of LRU, such as the second chance
  algorithm, which reads and clears the access bit; the dirty bit tells
  whether a victim must be written out at all.
\end{keypoint}

\section{Copy-on-write}
\label{sec:l08-cow}

The MMU translates addresses, tracks accesses and enforces protection.
These abilities are also useful for building other services and
optimizations.\source{P3L2, copy on write; OSTEP ch.~23; Silberschatz et
al.\ ch.~10.}  One of them avoids unnecessary copies when a process is
created.

A process created with \texttt{fork}\index{fork@\texttt{fork}} receives a copy of the
entire address space of its parent (Lesson~2).  Many of its pages, such as
the code, are static and never change, and a child often calls
\texttt{exec} soon after its creation, so most of the copying would be
wasted.

\begin{definition}[Copy-on-write]
  With \term{copy-on-write}[コピーオンライト]\index{COW|see{copy-on-write}}
  (COW) the operating system does not copy the address space when it
  creates a process.  It maps the virtual pages of the new process to the
  page frames of the original and write-protects these frames in both
  processes.  A write to such a page causes a page fault, and only then does
  the operating system copy the page.
\end{definition}

As long as both processes only read the shared pages, no memory is spent
and no time is spent on copying.  When one of them writes a page, the
protection fault lets the operating system recognize the COW page; it
copies this one page to a new page frame, maps the writing process's page
to the copy with write permission, and restarts the write
(\cref{fig:l08-cow}).  The cost of a copy is paid only for pages that are
actually modified.

\begin{figure}[htbp]
  \centering
  % Copy-on-write after fork: (a) parent and child map the same write-
% protected frames; (b) the child's write to page 1 faults, and the OS
% copies only that page.
% Usage: % Copy-on-write after fork: (a) parent and child map the same write-
% protected frames; (b) the child's write to page 1 faults, and the OS
% copies only that page.
% Usage: % Copy-on-write after fork: (a) parent and child map the same write-
% protected frames; (b) the child's write to page 1 faults, and the OS
% copies only that page.
% Usage: \input{figures/l08-cow}   (width ~118 mm)
\begin{tikzpicture}[x=1mm, y=1mm, font=\scriptsize\sffamily,
  >={Stealth[length=1.4mm]},
  ro/.style={draw, fill=ln-caution!12},
  rw/.style={draw, fill=ln-accent!22},
  oth/.style={draw, fill=ln-box},
  hd/.style={font=\scriptsize\sffamily\bfseries, anchor=base},
  capt/.style={font=\footnotesize\sffamily},
  lab/.style={font=\scriptsize\sffamily}]
  % \panel{x0}{child page-1 frame}{copy?}
  \def\panel#1#2{%
    \begin{scope}[shift={(#1,0)}]
      \node[hd] at (6,37.5) {\iflangja{親}{parent}};
      \node[hd] at (28,37.5) {\iflangja{メモリ}{memory}};
      \node[hd] at (50,37.5) {\iflangja{子}{child}};
      % frames 0..4, height 7, centres 3.5 + 7 i
      \foreach \i in {0,...,4} { \draw[oth] (22,7*\i) rectangle ++(12,7); }
      \foreach \i/\n in {3/0, 1/1, 4/2} {
        \draw[ro] (22,7*\i) rectangle ++(12,7) node[midway, lab] {\iflangja{ページ}{page} \n};
      }
      % parent pages 0..2 (centres 10, 18, 26)
      \foreach \j/\i in {0/3, 1/1, 2/4} {
        \draw[ro] (0,6+8*\j) rectangle ++(12,8) node[midway, lab] {\j};
        \draw[->] (12,10+8*\j) -- (22,3.5+7*\i);
      }
      #2
    \end{scope}}
  % (a) after fork: child maps the same frames
  \panel{0}{%
    \foreach \j/\i in {0/3, 1/1, 2/4} {
      \draw[ro] (44,6+8*\j) rectangle ++(12,8) node[midway, lab] {\j};
      \draw[->] (44,10+8*\j) -- (34,3.5+7*\i);
    }}
  \node[capt] at (28,-5) {\iflangja{(a) \texttt{fork} 直後}{(a) after \texttt{fork}}};
  % (b) after the child writes page 1
  \panel{62}{%
    \foreach \j/\i in {0/3, 2/4} {
      \draw[ro] (44,6+8*\j) rectangle ++(12,8) node[midway, lab] {\j};
      \draw[->] (44,10+8*\j) -- (34,3.5+7*\i);
    }
    \draw[rw] (44,14) rectangle ++(12,8) node[midway, lab] {1};
    \draw[rw] (22,0) rectangle ++(12,7) node[midway, lab] {\iflangja{コピー}{copy}};
    \draw[->, very thick, ln-accent] (44,18) -- (34,3.5);}
  \node[capt] at (90,-5) {\iflangja{(b) 子がページ 1 に書き込んだ後}{(b) after the child writes page 1}};
  % legend
  \draw[ro] (0,-12) rectangle ++(4,3);
  \node[lab, anchor=west] at (5,-10.5) {\iflangja{書き込み禁止}{write-protected}};
  \draw[rw] (32,-12) rectangle ++(4,3);
  \node[lab, anchor=west] at (37,-10.5) {\iflangja{書き込み可能}{writable}};
\end{tikzpicture}
   (width ~118 mm)
\begin{tikzpicture}[x=1mm, y=1mm, font=\scriptsize\sffamily,
  >={Stealth[length=1.4mm]},
  ro/.style={draw, fill=ln-caution!12},
  rw/.style={draw, fill=ln-accent!22},
  oth/.style={draw, fill=ln-box},
  hd/.style={font=\scriptsize\sffamily\bfseries, anchor=base},
  capt/.style={font=\footnotesize\sffamily},
  lab/.style={font=\scriptsize\sffamily}]
  % \panel{x0}{child page-1 frame}{copy?}
  \def\panel#1#2{%
    \begin{scope}[shift={(#1,0)}]
      \node[hd] at (6,37.5) {\iflangja{親}{parent}};
      \node[hd] at (28,37.5) {\iflangja{メモリ}{memory}};
      \node[hd] at (50,37.5) {\iflangja{子}{child}};
      % frames 0..4, height 7, centres 3.5 + 7 i
      \foreach \i in {0,...,4} { \draw[oth] (22,7*\i) rectangle ++(12,7); }
      \foreach \i/\n in {3/0, 1/1, 4/2} {
        \draw[ro] (22,7*\i) rectangle ++(12,7) node[midway, lab] {\iflangja{ページ}{page} \n};
      }
      % parent pages 0..2 (centres 10, 18, 26)
      \foreach \j/\i in {0/3, 1/1, 2/4} {
        \draw[ro] (0,6+8*\j) rectangle ++(12,8) node[midway, lab] {\j};
        \draw[->] (12,10+8*\j) -- (22,3.5+7*\i);
      }
      #2
    \end{scope}}
  % (a) after fork: child maps the same frames
  \panel{0}{%
    \foreach \j/\i in {0/3, 1/1, 2/4} {
      \draw[ro] (44,6+8*\j) rectangle ++(12,8) node[midway, lab] {\j};
      \draw[->] (44,10+8*\j) -- (34,3.5+7*\i);
    }}
  \node[capt] at (28,-5) {\iflangja{(a) \texttt{fork} 直後}{(a) after \texttt{fork}}};
  % (b) after the child writes page 1
  \panel{62}{%
    \foreach \j/\i in {0/3, 2/4} {
      \draw[ro] (44,6+8*\j) rectangle ++(12,8) node[midway, lab] {\j};
      \draw[->] (44,10+8*\j) -- (34,3.5+7*\i);
    }
    \draw[rw] (44,14) rectangle ++(12,8) node[midway, lab] {1};
    \draw[rw] (22,0) rectangle ++(12,7) node[midway, lab] {\iflangja{コピー}{copy}};
    \draw[->, very thick, ln-accent] (44,18) -- (34,3.5);}
  \node[capt] at (90,-5) {\iflangja{(b) 子がページ 1 に書き込んだ後}{(b) after the child writes page 1}};
  % legend
  \draw[ro] (0,-12) rectangle ++(4,3);
  \node[lab, anchor=west] at (5,-10.5) {\iflangja{書き込み禁止}{write-protected}};
  \draw[rw] (32,-12) rectangle ++(4,3);
  \node[lab, anchor=west] at (37,-10.5) {\iflangja{書き込み可能}{writable}};
\end{tikzpicture}
   (width ~118 mm)
\begin{tikzpicture}[x=1mm, y=1mm, font=\scriptsize\sffamily,
  >={Stealth[length=1.4mm]},
  ro/.style={draw, fill=ln-caution!12},
  rw/.style={draw, fill=ln-accent!22},
  oth/.style={draw, fill=ln-box},
  hd/.style={font=\scriptsize\sffamily\bfseries, anchor=base},
  capt/.style={font=\footnotesize\sffamily},
  lab/.style={font=\scriptsize\sffamily}]
  % \panel{x0}{child page-1 frame}{copy?}
  \def\panel#1#2{%
    \begin{scope}[shift={(#1,0)}]
      \node[hd] at (6,37.5) {\iflangja{親}{parent}};
      \node[hd] at (28,37.5) {\iflangja{メモリ}{memory}};
      \node[hd] at (50,37.5) {\iflangja{子}{child}};
      % frames 0..4, height 7, centres 3.5 + 7 i
      \foreach \i in {0,...,4} { \draw[oth] (22,7*\i) rectangle ++(12,7); }
      \foreach \i/\n in {3/0, 1/1, 4/2} {
        \draw[ro] (22,7*\i) rectangle ++(12,7) node[midway, lab] {\iflangja{ページ}{page} \n};
      }
      % parent pages 0..2 (centres 10, 18, 26)
      \foreach \j/\i in {0/3, 1/1, 2/4} {
        \draw[ro] (0,6+8*\j) rectangle ++(12,8) node[midway, lab] {\j};
        \draw[->] (12,10+8*\j) -- (22,3.5+7*\i);
      }
      #2
    \end{scope}}
  % (a) after fork: child maps the same frames
  \panel{0}{%
    \foreach \j/\i in {0/3, 1/1, 2/4} {
      \draw[ro] (44,6+8*\j) rectangle ++(12,8) node[midway, lab] {\j};
      \draw[->] (44,10+8*\j) -- (34,3.5+7*\i);
    }}
  \node[capt] at (28,-5) {\iflangja{(a) \texttt{fork} 直後}{(a) after \texttt{fork}}};
  % (b) after the child writes page 1
  \panel{62}{%
    \foreach \j/\i in {0/3, 2/4} {
      \draw[ro] (44,6+8*\j) rectangle ++(12,8) node[midway, lab] {\j};
      \draw[->] (44,10+8*\j) -- (34,3.5+7*\i);
    }
    \draw[rw] (44,14) rectangle ++(12,8) node[midway, lab] {1};
    \draw[rw] (22,0) rectangle ++(12,7) node[midway, lab] {\iflangja{コピー}{copy}};
    \draw[->, very thick, ln-accent] (44,18) -- (34,3.5);}
  \node[capt] at (90,-5) {\iflangja{(b) 子がページ 1 に書き込んだ後}{(b) after the child writes page 1}};
  % legend
  \draw[ro] (0,-12) rectangle ++(4,3);
  \node[lab, anchor=west] at (5,-10.5) {\iflangja{書き込み禁止}{write-protected}};
  \draw[rw] (32,-12) rectangle ++(4,3);
  \node[lab, anchor=west] at (37,-10.5) {\iflangja{書き込み可能}{writable}};
\end{tikzpicture}

  \caption{Copy-on-write.  (a) After \texttt{fork}, parent and child map the
    same write-protected page frames.  (b) The child's write to page~1
    faults, and the operating system copies only that page.}
  \label{fig:l08-cow}
\end{figure}

\begin{example}
  A process with an address space of 25\,600 pages of \SI{4}{\kibi\byte}
  (\SI{100}{\mebi\byte}) forks, and the child writes to 40 pages of its
  stack and data before it calls \texttt{exec}.  Copy-on-write copies
  40 pages instead of 25\,600, i.e., \SI{0.16}{\percent} of the address
  space.
\end{example}

\section{Checkpointing}
\label{sec:l08-checkpointing}

The same hardware support helps with failure
management.\source{P3L2, failure management: checkpointing.}

\begin{definition}[Checkpointing]
  A failure and recovery management technique that periodically saves the
  state of a process is called \term{checkpointing}[チェックポインティング].
  After a failure, the process is restarted from its most recent checkpoint
  rather than from the beginning.
\end{definition}

Failures may be unavoidable, but with checkpoints recovery is much faster.
The simplest approach pauses the process and copies its entire state.  The
pause is long for a large address space, and the approach is expensive if
it is repeated.  The MMU enables a better approach.
\begin{itemize}
  \item The operating system write-protects all pages of the process and
    copies the entire state once, while the process continues to run.
  \item From then on it records which pages the process modifies, through
    the write faults or the dirty bits, and later checkpoints copy only
    these dirtied pages.  Such \emph{incremental} checkpoints store the
    differences to the previous checkpoint.
  \item To restore the process, its state is rebuilt from the full copy and
    the subsequent differences; the differences can also be merged in the
    background.
\end{itemize}

Checkpointing is the basis of further services.
\begin{description}
  \item[Debugging] With \emph{rewind-replay}, a program is restarted
    (rewound) from a checkpoint and re-executed.  If the error does not
    reappear, successively older checkpoints are used until it is found.
  \item[Migration] With \term{process migration}[プロセスマイグレーション]
    the checkpoint of a process is transferred to another machine, where
    the process continues.  Migration serves disaster recovery and
    consolidation of processes onto fewer machines.
\end{description}
To keep the pause of a migration short, the checkpoints are repeated in a
fast loop: the operating system copies all pages while the process runs,
then the pages dirtied during that copy, then those dirtied during the
second copy, and so on, until so few pages remain that a final
pause-and-copy is acceptable---or until it is unavoidable because the
process dirties pages as fast as they are copied.

\needspace{14\baselineskip}
\begin{example}
  A process of 50\,000 pages is migrated over a connection that copies
  25\,000 pages per second, while the process dirties 2000 distinct pages
  per second.  Pausing the process and copying everything takes
  \SI{2}{\second}.  With repeated rounds, each of which copies the pages
  dirtied during the previous one, every round copies \SI{8}{\percent} of
  the pages of the round before, and after three rounds the final pause
  lasts about \SI{1}{\milli\second}:
  \begin{center}
    \small
    \begin{tabular}{@{}lrrr@{}}
      \toprule
      Round & Pages copied & Duration & Pages dirtied meanwhile \\
      \midrule
      1 (process runs)   & 50\,000 & \SI{2}{\second}             & 4000 \\
      2 (process runs)   & 4000    & \SI{160}{\milli\second}     & 320 \\
      3 (process runs)   & 320     & \SI{12.8}{\milli\second}    & 26 \\
      final (paused)     & 26      & \SI{1.04}{\milli\second}    & --- \\
      \bottomrule
    \end{tabular}
  \end{center}
\end{example}

\needspace{14\baselineskip}

\begin{keypoint}[Lesson summary]
  Virtual memory decouples address spaces from physical memory; the
  operating system allocates memory and arbitrates every access, using
  paging or segmentation and the MMU, which translates addresses and raises
  faults.  A page table maps VPNs to PFNs (\cref{eq:l08-translation}), and
  its entries carry valid, dirty, access and protection bits; a page fault
  hands an access that the MMU cannot complete to the operating system.
  Flat page tables are too large (\cref{eq:l08-pt-size}), so they are
  multi-level or inverted and hashed, and the TLB keeps translation fast
  (\cref{eq:l08-eat}).  Large pages shrink page tables and raise TLB
  coverage at the cost of internal fragmentation.  Allocators fight external
  fragmentation; Linux uses a buddy allocator for page frames and a slab
  allocator for kernel objects.  Demand paging swaps pages in on a page fault,
  and a page-out daemon swaps pages out, choosing them with LRU
  approximations such as the second chance algorithm.  The MMU also enables
  copy-on-write and incremental checkpointing, which in turn supports
  debugging and process migration.
\end{keypoint}

\end{document}
